\chapter{Isotropic and scattering calculi}\label{ch:isotropic}
\providecommand{\Gs}[1]{\Gamma^{#1}}
\providecommand{\Gph}[1]{\Gamma^{#1}_{\mathrm{ph}}}
\providecommand{\Psiisoph}[1]{\Psi^{#1}_{\mathrm{iso,ph}}}
\providecommand{\Hiso}[1]{H^{#1}_{\mathrm{iso}}}
\providecommand{\Gres}{G^{-\infty}_{\mathrm{iso}}}
\providecommand{\Ampl}[1]{\mathcal{A}^{#1}}
\providecommand{\rc}{\operatorname{RC}}
\providecommand{\Hsph}[1]{\mathbb{S}^{#1}_{+}}
\providecommand{\OpW}{\operatorname{Op}^{\mathrm{W}}}

In this chapter two \emph{global} calculi of pseudodifferential operators on
$\R^n$ are constructed and analysed. In the \emph{isotropic} calculus the
symbols are required to satisfy estimates in which the space variables $x$ and
the dual variables $\xi$ play completely symmetric roles: each derivative, in
$x$ or in $\xi$, improves the decay by a factor of $(1+|x|+|\xi|)^{-1}$. In
the \emph{scattering} calculus derivatives in $x$ gain powers of
$\ang{x}^{-1}$ and derivatives in $\xi$ gain powers of $\ang{\xi}^{-1}$,
separately. Both calculi are subalgebras of (weighted versions of) the
traditional calculus $\Psido{m}(\R^n)$ of \cref{ch:euclidean}, and we use the
results proved there freely. The isotropic calculus is especially simple
because its residual algebra $\Psiiso{-\infty}(\R^n)$ consists of the
operators with kernels in $\cS(\R^{2n})$; these operators are compact (indeed
Hilbert--Schmidt) on $L^2(\R^n)$. This residual algebra, the associated group
$\Gres(\R^n)$ of invertible perturbations of the identity, and the
approximation of residual operators by finite rank operators built from the
eigenfunctions of the harmonic oscillator, are the basic objects for the
discussion of traces and determinants in \cref{ch:determinants}. Thus this
chapter is preparatory: the constructions (adjoints, composition,
parametrices, generalized inverses, Sobolev spaces, $L^2$ estimates, complex
powers) reappear later for operators on manifolds, and here they can be
carried out without the complications of coordinate changes.

Unlike the original notes, we work with general (not necessarily
polyhomogeneous) symbols by default and indicate the polyhomogeneous
subclasses separately.

\medskip\noindent\textbf{Notation.} Points of $\R^{2n}=\R^n_x\times\R^n_\xi$
are written $w=(x,\xi)$ and $\ang{w}=\ang{(x,\xi)}=(1+|x|^2+|\xi|^2)^{1/2}$.
We repeatedly use the following elementary inequalities. For $v,w\in\R^N$ and
$s\in\R$ (Peetre's inequality)
\begin{equation}\label{isotropic:eq:peetre}
  \ang{w+v}^s\le 2^{|s|/2}\ang{w}^s\ang{v}^{|s|},
\end{equation}
which follows from $1+|w+v|^2\le 2(1+|w|^2)(1+|v|^2)$, applied to $(w,v)$ if
$s\ge0$ and to $(w+v,-v)$ if $s<0$. For $z\in\R^p$, $\xi\in\R^n$ and
$\theta\in[0,1]$,
\begin{equation}\label{isotropic:eq:split}
  \max(\ang{z},\ang{\xi})\le\ang{(z,\xi)}\le\ang{z}\ang{\xi},\qquad
  \ang{(z,\xi)}\ge\ang{z}^\theta\ang{\xi}^{1-\theta},
\end{equation}
the last by the arithmetic--geometric mean inequality,
$\ang{z}^{2\theta}\ang{\xi}^{2(1-\theta)}\le\theta\ang{z}^2+(1-\theta)\ang{\xi}^2\le\ang{(z,\xi)}^2$.
We write $m_+=\max(m,0)$.

%======================================================================
\section{Symbol classes}\label{isotropic:sec:symbols}

\subsection{Isotropic symbols}

\begin{definition}\label{isotropic:def:Gamma}
For $m\in\R$ and $N\in\N$ the space of \emph{isotropic symbols} of order $m$ is
\begin{equation*}
  \Gs{m}(\R^N)=\bigl\{a\in\CI(\R^N);\ |\pa^\alpha a(w)|\le C_\alpha\ang{w}^{m-|\alpha|}\ \forall\,\alpha\in\N_0^N\bigr\}.
\end{equation*}
It is a Fréchet space with the norms
$\norm{a}_{m,k}=\max_{|\alpha|\le k}\sup_w\ang{w}^{-m+|\alpha|}|\pa^\alpha a(w)|$.
We set $\Gs{-\infty}(\R^N)=\bigcap_m\Gs{m}(\R^N)=\cS(\R^N)$ and
$\Gs{\infty}(\R^N)=\bigcup_m\Gs{m}(\R^N)$. Vector- and matrix-valued symbols are
defined in the same way.
\end{definition}

\begin{lemma}\label{isotropic:lem:Gamma-basic}
\begin{enumerate}[label=\textup{(\roman*)}]
\item $\Gs{m}\cdot\Gs{m'}\subset\Gs{m+m'}$, $\pa^\alpha:\Gs{m}\to\Gs{m-|\alpha|}$,
  $\Gs{m}\subset\Gs{m'}$ for $m\le m'$, and $\ang{w}^s\in\Gs{s}$ for all $s\in\R$.
\item If $a\in\Gs{0}(\R^N;\C^M)$ and $F$ is smooth on a neighbourhood of the
  (compact) closure of $a(\R^N)$, then $F\circ a\in\Gs{0}(\R^N)$.
\item If $L$ is a linear isomorphism of $\R^N$ then $a\mapsto a\circ L$ is an
  isomorphism of $\Gs{m}(\R^N)$.
\item $\cS(\R^N)$ is dense in $\Gs{m}(\R^N)$ in the topology of $\Gs{m'}(\R^N)$
  for any $m'>m$.
\end{enumerate}
\end{lemma}
\begin{proof}
(i) follows from Leibniz' formula and, for $\ang{w}^s$, from the fact that
$\pa^\alpha\ang{w}^s$ is a sum of terms $p(w)\ang{w}^{s-2k}$ with $p$ a
polynomial of degree at most $2k-|\alpha|$. (ii) By the chain rule
$\pa^\alpha(F\circ a)$ is a sum of terms
$(\pa^kF)(a)\,\pa^{\alpha_1}a\cdots\pa^{\alpha_k}a$ with
$\alpha_1+\dots+\alpha_k=\alpha$, $|\alpha_i|\ge1$, each bounded by
$C\ang{w}^{-|\alpha|}$. (iii) holds since $c\ang{w}\le\ang{Lw}\le C\ang{w}$.
(iv) Let $\chi\in\CIc(\R^N)$ equal $1$ near $0$ and $a_k=\chi(w/k)a$. The
functions $\chi(w/k)$ are bounded in $\Gs{0}$ uniformly in $k\ge1$ and
$1-\chi(w/k)$ vanishes for $|w|\le ck$, where
$\ang{w}^{m-m'}\le Ck^{m-m'}$; hence $\norm{a-a_k}_{m',j}\le C_jk^{m-m'}\to0$.
\end{proof}

\begin{lemma}[Asymptotic completeness]\label{isotropic:lem:asymp-Gamma}
If $a_j\in\Gs{m-j}(\R^N)$ for $j\in\N_0$, there exists $a\in\Gs{m}(\R^N)$ such that
\begin{equation*}
  a-\sum_{j<k}a_j\in\Gs{m-k}(\R^N)\quad\forall\,k\in\N_0;
\end{equation*}
$a$ is unique modulo $\cS(\R^N)$ and we write $a\sim\sum_ja_j$.
\end{lemma}
\begin{proof}
Choose $\psi\in\CI(\R^N)$ with $\psi(w)=0$ for $|w|\le1$, $\psi(w)=1$ for
$|w|\ge2$, and set $\psi_\eps(w)=\psi(\eps w)$ for $0<\eps\le1$. Then
$\pa^\alpha\psi_\eps=\eps^{|\alpha|}(\pa^\alpha\psi)(\eps w)$ is, for
$\alpha\ne0$, supported in $1/\eps\le|w|\le2/\eps$, so
$|\pa^\alpha\psi_\eps|\le C_\alpha\ang{w}^{-|\alpha|}$ uniformly in $\eps$;
moreover $\ang{w}\ge1/\eps$ on $\supp\psi_\eps$. Hence for $k<j$ and
$\ell\in\N_0$
\begin{equation*}
  \norm{\psi_\eps a_j}_{m-k,\ell}\le C_{j,\ell}\sup_{\supp\psi_\eps}\ang{w}^{k-j}\le C_{j,\ell}\,\eps .
\end{equation*}
Choose $\eps_j\downarrow0$ with $C_{j,j}\eps_j\le 2^{-j}$ (we may take
$C_{j,\ell}$ increasing in $\ell$ and valid for all $k<j$). Then for each $k$
and $\ell$ the series $\sum_{j>k}\psi_{\eps_j}a_j$ converges in the norm
$\norm{\cdot}_{m-k,\ell}$, since all but finitely many terms are bounded by
$2^{-j}$. Thus $a=\sum_j\psi_{\eps_j}a_j$ converges in $\Gs{m}$ and
\begin{equation*}
  a-\sum_{j<k}a_j=\sum_{j<k}(\psi_{\eps_j}-1)a_j+\psi_{\eps_k}a_k+\sum_{j>k}\psi_{\eps_j}a_j
  \in\CIc+\Gs{m-k}+\Gs{m-k}.
\end{equation*}
Uniqueness holds because $\bigcap_k\Gs{m-k}=\cS(\R^N)$.
\end{proof}
Sums with steps of two, $a\sim\sum_ja_j$ with $a_j\in\Gs{m-2j}$, are covered
by inserting zero terms.

\subsection{Scattering symbols}

\begin{definition}\label{isotropic:def:Ssc}
For $m,l\in\R$ the space $S^{m,l}(\R^p;\R^n)$ of \emph{scattering symbols}
consists of the $a\in\CI(\R^p_z\times\R^n_\xi)$ such that
\begin{equation*}
  |D_z^\alpha D_\xi^\beta a(z,\xi)|\le C_{\alpha\beta}\ang{z}^{l-|\alpha|}\ang{\xi}^{m-|\beta|}\quad\forall\,\alpha,\beta.
\end{equation*}
Here $m$ is the order in the fibre variables and $l$ the order in the base
variables.
\end{definition}

Clearly $S^{m,l}\cdot S^{m',l'}\subset S^{m+m',l+l'}$,
$D_z:S^{m,l}\to S^{m,l-1}$, $D_\xi:S^{m,l}\to S^{m-1,l}$,
$S^{m,l}(\R^p;\R^n)\subset\ang{z}^lS^m(\R^p;\R^n)$, and
$\bigcap_{m,l}S^{m,l}(\R^n;\R^n)=\cS(\R^{2n})$.

\begin{lemma}\label{isotropic:lem:asymp-sc}
If $a_j\in S^{m-j,l-j}(\R^p;\R^n)$, $j\in\N_0$, there exists $a\in S^{m,l}(\R^p;\R^n)$
with $a-\sum_{j<k}a_j\in S^{m-k,l-k}(\R^p;\R^n)$ for all $k$; $a$ is unique
modulo $\cS(\R^{p+n})$.
\end{lemma}
\begin{proof}
Repeat the proof of \cref{isotropic:lem:asymp-Gamma} with
$\psi_\eps(z,\xi)=1-(1-\psi(\eps z))(1-\psi(\eps\xi))$, which is bounded in
$S^{0,0}$ uniformly in $\eps$, vanishes where $|z|\le1/\eps$ and
$|\xi|\le1/\eps$ (so $1-\psi_\eps$ has compact support), and satisfies
$\ang{z}^{-1}\ang{\xi}^{-1}\le\eps$ on its support. Thus for $k<j$ the
$S^{m-k,l-k}$ norms of $\psi_\eps a_j$ are bounded by
$C\sup_{\supp\psi_\eps}(\ang{z}\ang{\xi})^{k-j}\le C\eps$.
\end{proof}

The two classes are related as follows; this replaces the corresponding lemma
of the original notes, where the power of $1+|x|^2$ is misprinted.

\begin{lemma}\label{isotropic:lem:Gamma-in-S}
For $z\in\R^p$, $\xi\in\R^n$,
\begin{equation*}
  \Gs{m}(\R^{p+n})\subset
  \begin{cases}S^{m,m}(\R^p;\R^n),& m\ge0,\\
   \bigcap_{m\le t\le0}S^{m-t,t}(\R^p;\R^n),& m\le0.\end{cases}
\end{equation*}
In particular $\Gs{m}(\R^{p+n})\subset\ang{z}^{m_+}S^m(\R^p;\R^n)$.
\end{lemma}
\begin{proof}
By \eqref{isotropic:eq:split},
$|\pa_z^\alpha\pa_\xi^\beta a|\le C\ang{(z,\xi)}^{m}\ang{z}^{-|\alpha|}\ang{\xi}^{-|\beta|}$.
If $m\ge0$, $\ang{(z,\xi)}^m\le\ang{z}^m\ang{\xi}^m$. If $m<0$ and $m\le t\le0$,
take $\theta=t/m\in[0,1]$ in \eqref{isotropic:eq:split} to get
$\ang{(z,\xi)}^m\le\ang{z}^{t}\ang{\xi}^{m-t}$.
\end{proof}

\subsection{Polyhomogeneous symbols and radial compactification}

A function $a\in\CI(\R^N)$ is \emph{homogeneous of degree $\mu$ outside the
unit ball} if $a(tw)=t^\mu a(w)$ for $|w|\ge1$, $t\ge1$. Such a function lies in
$\Gs{\mu}(\R^N)$, since on $|w|\ge1$ its derivatives of order $k$ are
homogeneous of degree $\mu-k$.

\begin{definition}\label{isotropic:def:Gph}
The space $\Gph{m}(\R^N)$ of \emph{polyhomogeneous} (classical) isotropic
symbols consists of those $a\in\Gs{m}(\R^N)$ for which there are $a_{m-j}$,
homogeneous of degree $m-j$ outside the unit ball, with $a\sim\sum_ja_{m-j}$.
The \emph{principal symbol} of $a$ is the restriction of $a_m$ to the sphere
$\Sphere^{N-1}$.
\end{definition}

The polyhomogeneous classes are closed under multiplication and
differentiation ($\pa_{w_j}:\Gph{m}\to\Gph{m-1}$), contain $\ang{w}^s\in\Gph{s}$
(expand $\ang{w}^s=|w|^s(1+|w|^{-2})^{s/2}$ binomially for $|w|\ge2$ and cut
the terms off near $0$), and if $a_j\in\Gph{m-j}$ and $a\sim\sum_ja_j$ then
$a\in\Gph{m}$, as one sees by collecting the finitely many terms of each
degree of homogeneity in the expansions of the $a_j$.

The \emph{radial compactification} is the map
\begin{equation*}
  \rc:\R^N\ni w\longmapsto\bigl(\ang{w}^{-1},w\ang{w}^{-1}\bigr)\in
  \Hsph{N}=\{X=(X_0,X')\in\R^{1+N};\ |X|=1,\ X_0\ge0\},
\end{equation*}
a diffeomorphism of $\R^N$ onto the interior of the closed half-sphere
$\Hsph{N}$, a compact manifold with boundary
$\pa\Hsph{N}=\{X_0=0\}\cong\Sphere^{N-1}$. The function $\rho=X_0$ is a
boundary defining function and $\rho\circ\rc=\ang{w}^{-1}$. We denote by
$\dCI(\Hsph{N})$ the smooth functions vanishing to infinite order at the
boundary. The space $\mathcal L(\Hsph{N})=\bigcup_{k\in\N_0}\rho^{-k}\CI(\Hsph{N})$
is the space of \emph{Laurent functions} on the half-sphere; more generally
$\rho^{\mu}\CI(\Hsph{N})$, $\mu\in\C$, is the space of products $\rho^\mu v$ with
$v$ smooth, regarded as functions on the interior.

\begin{proposition}\label{isotropic:prop:rc}
\begin{enumerate}[label=\textup{(\roman*)}]
\item For every $m\in\R$, $\rc^*\bigl(\rho^{-m}\CI(\Hsph{N})\bigr)=\Gph{m}(\R^N)$, and the
  principal symbol of $\rc^*(\rho^{-m}F)$ is $F|_{\pa\Hsph{N}}$.
\item $\rc^*\dCI(\Hsph{N})=\cS(\R^N)$.
\item If $f\in\Gs{-k}(\R^N)$, $k\in\N$, then $f\circ\rc^{-1}$ extends to a
  $\mathcal C^{k-1}$ function on $\Hsph{N}$ vanishing at the boundary, and
  $\Gs{-k}(\R^N)\ni f\mapsto f\circ\rc^{-1}\in\mathcal C^{k-1}(\Hsph{N})$ is continuous.
\end{enumerate}
\end{proposition}
\begin{proof}[Sketch of proof]
Near the boundary use the coordinates $\rho\in[0,1)$ and
$\theta=X'/|X'|\in\Sphere^{N-1}$; in terms of $w$, $\rho=\ang{w}^{-1}$ and
$\theta=w/|w|$, so $|w|=(1-\rho^2)^{1/2}/\rho$. In polar coordinates
$\pa_{w_j}=\theta_j\pa_r+r^{-1}\Omega_j$, where $\Omega_j$ is the vector field
on $\Sphere^{N-1}$ obtained by projecting $e_j$ tangentially, and
$\pa_r=-r\rho^3\pa_\rho$. Hence
\begin{equation}\label{isotropic:eq:b-fields}
  V_j:=\ang{w}\pa_{w_j}=-\theta_j(1-\rho^2)^{1/2}\rho\pa_\rho+(1-\rho^2)^{-1/2}\Omega_j .
\end{equation}
These are smooth vector fields on $\Hsph{N}$ tangent to the boundary, and
since $\sum_j\theta_j\Omega_j=0$ they span, over $\CI$, the fields $\rho\pa_\rho$
and $\Omega_j$, that is all smooth vector fields tangent to the boundary. An
induction shows that $\ang{w}^{|\alpha|}\pa^\alpha$ is a combination, with
coefficients in $\Gs{0}$, of products of at most $|\alpha|$ of the $V_j$.
Consequently $\rc^*\CI(\Hsph{N})\subset\Gs{0}$. Conversely, a derivative of
order $j$ of $f\circ\rc^{-1}$ in $(\rho,\theta)$ is $\rho^{-j}$ times a
combination of at most $j$-fold products of the $V_j$ applied to $f$; this
gives (iii), since such products applied to $f\in\Gs{-k}$ are
$O(\ang{w}^{-k})=O(\rho^k)$, and also (ii), since for $f\in\cS$ these are
$O(\rho^K)$ for every $K$ (the converse inclusion in (ii) follows from the
first argument with $\rho^{-K}F$, $F\in\dCI$).

For (i), if $F\in\CI(\Hsph{N})$, Taylor's formula in $\rho$ gives
$F=\sum_{j<K}\rho^jF_j(\theta)+\rho^KR_K$ near the boundary. Since
$\rho=|w|^{-1}(1+|w|^{-2})^{-1/2}$ for $|w|\ge1$, each $\rho^{j-m}F_j(\theta)$
pulls back to an asymptotic sum of functions homogeneous of degrees
$m-j-2i$, $i\in\N_0$, while $\rho^{K-m}R_K$ pulls back into $\Gs{m-K}$. Thus
$\rc^*(\rho^{-m}F)\in\Gph{m}$, with leading term $|w|^mF(0,\theta)$. Conversely,
a function homogeneous of degree $m-j$ outside the unit ball is
$|w|^{m-j}f(\theta)=\rho^{j-m}(1-\rho^2)^{(m-j)/2}f(\theta)$, which lies in
$\rho^{-m}\rho^j\CI$ near the boundary. Given $a\sim\sum_ja_{m-j}$, Borel's
lemma gives $F\in\CI(\Hsph{N})$ with $\rho^{-m}F$ having the same expansion,
and then $a-\rc^*(\rho^{-m}F)\in\cS(\R^N)=\rc^*\dCI$ by (ii).
\end{proof}

Thus $\Gph{\Z}(\R^N)=\bigcup_{k\in\Z}\Gph{k}(\R^N)$ is identified with the
Laurent functions $\mathcal L(\Hsph{N})$ and $\Gph{0}(\R^N)$ with
$\CI(\Hsph{N})$.

%======================================================================
\section{Isotropic operators}\label{isotropic:sec:iso-ops}

Recall the left quantization
$\Op(a)u(x)=\int e^{i(x-y)\cdot\xi}a(x,\xi)u(y)\,dy\,\dbar\xi$ and, for an
amplitude $a(x,y,\xi)$, the same formula with $a(x,y,\xi)$.

\begin{definition}\label{isotropic:def:Psiiso}
For $m\in\R$ the space of \emph{isotropic pseudodifferential operators} of
order $m$ on $\R^n$ is
\begin{equation*}
  \Psiiso{m}(\R^n)=\{\Op(a);\ a\in\Gs{m}(\R^{2n})\},\qquad
  \Psiisoph{m}(\R^n)=\{\Op(a);\ a\in\Gph{m}(\R^{2n})\}.
\end{equation*}
We set $\Psiiso{-\infty}=\bigcap_m\Psiiso{m}$ and $\Psiiso{\infty}=\bigcup_m\Psiiso{m}$.
\end{definition}

By \cref{isotropic:lem:Gamma-in-S}, if $a\in\Gs{m}(\R^{2n})$ then
$a=\ang{x}^{m_+}b$ with $b\in S^m(\R^n;\R^n)$, so
$\Op(a)=\ang{x}^{m_+}\Op(b)\in\ang{x}^{m_+}\Psido{m}(\R^n)$. In particular
isotropic operators act continuously on $\cS(\R^n)$ and on $\cSp(\R^n)$, by the
results of \cref{ch:euclidean}. The kernel of $\Op(a)$ is
$K(x,y)=\int e^{i(x-y)\cdot\xi}a(x,\xi)\,\dbar\xi$, i.e.\ the partial inverse
Fourier transform of $a$ in $\xi$ evaluated at $(x,x-y)$; hence $a$ is
determined by $\Op(a)$ and we call $\sigma_L(\Op(a))=a$ the \emph{(full) left
symbol}. Since partial Fourier transformation and linear changes of variables
preserve $\cS$:

\begin{lemma}\label{isotropic:lem:residual-kernels}
An operator $A$ is in $\Psiiso{-\infty}(\R^n)$ if and only if it has Schwartz
kernel in $\cS(\R^{2n})$, and $\Psiiso{-\infty}(\R^n)\ni A\mapsto K_A\in\cS(\R^{2n})$
is a bijection.
\end{lemma}

We topologize $\Psiiso{-\infty}(\R^n)$ by this identification.

\begin{example}\label{isotropic:ex:differential}
Since $\Op(x^\beta\xi^\alpha)=x^\beta D^\alpha$, every differential operator
with polynomial coefficients,
$P=\sum_{|\alpha|+|\beta|\le m}c_{\alpha\beta}x^\beta D^\alpha$, lies in
$\Psiisoph{m}(\R^n)$. Conversely, a differential operator lies in
$\Psiiso{\infty}(\R^n)$ only if its coefficients are polynomials (a smooth
function with all derivatives of order $k$ bounded by $C\ang{x}^{m-k}$ for all
$k$ is a polynomial). The basic example is the \emph{harmonic oscillator}
$H=\sum_j(D_j^2+x_j^2)$, with symbol $|x|^2+|\xi|^2$, studied in
\cref{isotropic:sec:harmonic}.
\end{example}

\subsection{Amplitudes and reduction}

As in \cref{ch:euclidean}, adjoints and products are obtained by reducing
operators defined by amplitudes to left-reduced form. Note that a left symbol
$a(x,\xi)\in\Gs{m}(\R^{2n})$, regarded as a function of $(x,y,\xi)$, is
\emph{not} in general in $\Gs{m}(\R^{3n})$, so we use a hybrid class of
amplitudes.

\begin{definition}\label{isotropic:def:ampl}
For $m\in\R$, $\Ampl{m}(\R^{3n})$ is the space of $a\in\CI(\R^{3n}_{x,y,\xi})$
for which there is $w\ge0$ such that for all multi-indices
$\alpha'\le\alpha$, $\beta'\le\beta$, $\gamma'\le\gamma$
\begin{equation}\label{isotropic:eq:ampl}
  |D_x^\alpha D_y^\beta D_\xi^\gamma a|\le
  C\ang{(x,\xi)}^{m-k'}\ang{x-y}^{w+k'}\ang{\xi}^{-|\gamma-\gamma'|},\qquad k'=|\alpha'|+|\beta'|+|\gamma'| .
\end{equation}
\end{definition}

Thus any subset of the derivatives may be used to improve the decay in
$(x,\xi)$, at the cost of a corresponding growth in $|x-y|$, while the
remaining $\xi$-derivatives still give decay in $\xi$.

\begin{lemma}\label{isotropic:lem:ampl-examples}
\begin{enumerate}[label=\textup{(\roman*)}]
\item If $a\in\Gs{m}(\R^{2n})$ and $b\in\Gs{m'}(\R^{2n})$ then
  $a(x,\xi)b(y,\xi)\in\Ampl{m+m'}(\R^{3n})$.
\item If $a\in\Gs{m}(\R^{2n})$ and $\lambda\in[0,1]$ then
  $a(x+\lambda(y-x),\xi)\in\Ampl{m}(\R^{3n})$.
\item If $a\in\Ampl{m}(\R^{3n})$ then
  $|D_x^\alpha D_y^\beta D_\xi^\gamma a|\le C\ang{x}^{m_+}\ang{x-y}^w\ang{\xi}^{m-|\gamma|}$.
\end{enumerate}
\end{lemma}
\begin{proof}
(i) By Leibniz' formula $D_x^\alpha D_y^\beta D_\xi^\gamma(ab)$ is a sum of
terms $(D_x^\alpha D_\xi^{\gamma_1}a)(x,\xi)(D_y^\beta D_\xi^{\gamma_2}b)(y,\xi)$
with $\gamma_1+\gamma_2=\gamma$. Given $\gamma'\le\gamma$, choose $\gamma_i'\le\gamma_i$
with $|\gamma_1'|+|\gamma_2'|=|\gamma'|$. Using
$\ang{(x,\xi)}^{-1}\le\ang{\xi}^{-1}$ we get
$|D_x^\alpha D_\xi^{\gamma_1}a|\le C\ang{(x,\xi)}^{m-|\alpha'|-|\gamma_1'|}\ang{\xi}^{-|\gamma_1-\gamma_1'|}$,
and similarly, followed by \eqref{isotropic:eq:peetre} with $v=(y-x,0)$,
\begin{equation*}
  |D_y^\beta D_\xi^{\gamma_2}b(y,\xi)|\le C\ang{(y,\xi)}^{m'-|\beta'|-|\gamma_2'|}\ang{\xi}^{-|\gamma_2-\gamma_2'|}
  \le C'\ang{(x,\xi)}^{m'-|\beta'|-|\gamma_2'|}\ang{x-y}^{|m'|+|\beta'|+|\gamma_2'|}\ang{\xi}^{-|\gamma_2-\gamma_2'|}.
\end{equation*}
Multiplying gives \eqref{isotropic:eq:ampl} with $w=|m'|$. (ii) With
$\bar x=x+\lambda(y-x)$ the derivative is a constant times
$(\pa^{\alpha+\beta}_xD_\xi^\gamma a)(\bar x,\xi)$, bounded by
$C\ang{(\bar x,\xi)}^{m-k'}\ang{\xi}^{-|\gamma-\gamma'|}$, and
\eqref{isotropic:eq:peetre} gives
$\ang{(\bar x,\xi)}^{m-k'}\le C\ang{(x,\xi)}^{m-k'}\ang{x-y}^{|m|+k'}$.
(iii) Take $\alpha'=\beta'=\gamma'=0$ and use
$\ang{(x,\xi)}^m\le\ang{x}^{m_+}\ang{\xi}^m$.
\end{proof}

By (iii), $\ang{x}^{-m_+}a$ is an amplitude of the traditional calculus of
\cref{ch:euclidean} (of order $m$ with weight $\ang{x-y}^w$), so
$\Op(a)=\ang{x}^{m_+}\Op(\ang{x}^{-m_+}a)\in\ang{x}^{m_+}\Psido{m}(\R^n)$ is a
well-defined continuous operator on $\cS(\R^n)$ and on $\cSp(\R^n)$.

\begin{proposition}[Reduction]\label{isotropic:prop:reduction}
If $a\in\Ampl{m}(\R^{3n})$ then $\Op(a)\in\Psiiso{m}(\R^n)$ and
\begin{equation}\label{isotropic:eq:left-reduction}
  \sigma_L(\Op(a))\sim\sum_\alpha\frac{i^{|\alpha|}}{\alpha!}\bigl(D_\xi^\alpha D_y^\alpha a\bigr)(x,x,\xi),
\end{equation}
where the term of index $\alpha$ lies in $\Gs{m-2|\alpha|}(\R^{2n})$.
\end{proposition}
\begin{proof}
Let $c_\alpha(x,\xi)=\frac{i^{|\alpha|}}{\alpha!}(D_\xi^\alpha D_y^\alpha a)(x,x,\xi)$.
A derivative $D_x^\mu D_\xi^\gamma c_\alpha$ is a sum of values at $y=x$ of
derivatives of $a$ containing $|\alpha|$ derivatives in $y$, $|\alpha|+|\gamma|$
in $\xi$ and $|\mu|$ in $x$ or $y$; applying \eqref{isotropic:eq:ampl} with all
derivatives used for decay, and $\ang{x-y}=1$ at $y=x$, shows
$c_\alpha\in\Gs{m-2|\alpha|}(\R^{2n})$.

Taylor's formula in $y$ around $x$ gives
\begin{equation*}
  a(x,y,\xi)=\sum_{|\alpha|<N}\frac{(y-x)^\alpha}{\alpha!}(\pa_y^\alpha a)(x,x,\xi)
  +\sum_{|\alpha|=N}\frac{N(y-x)^\alpha}{\alpha!}\int_0^1(1-t)^{N-1}(\pa_y^\alpha a)(x,x+t(y-x),\xi)\,dt .
\end{equation*}
Since $(y-x)^\alpha e^{i(x-y)\cdot\xi}=(-D_\xi)^\alpha e^{i(x-y)\cdot\xi}$,
integration by parts in $\xi$ (legitimate for these oscillatory integrals, as
in \cref{ch:euclidean}) gives
\begin{equation}\label{isotropic:eq:taylor-remainder}
  \Op(a)=\sum_{|\alpha|<N}\Op(c_\alpha)+\Op(r_N),\quad
  r_N(x,y,\xi)=\sum_{|\alpha|=N}\frac{N i^N}{\alpha!}\int_0^1(1-t)^{N-1}
  \bigl(D_\xi^\alpha D_y^\alpha a\bigr)(x,x+t(y-x),\xi)\,dt .
\end{equation}
We estimate $r_N$. A derivative $D_x^\mu D_y^\nu D_\xi^\gamma$ of the
integrand is, by the chain rule, a combination with bounded coefficients of
derivatives of $a$ evaluated at $(x,x+t(y-x),\xi)$ which contain the $N$
derivatives $D_y^\alpha$, the $N+|\gamma|$ derivatives $D_\xi^{\alpha+\gamma}$
and further $x$- and $y$-derivatives. Applying \eqref{isotropic:eq:ampl} using
exactly the $2N$ derivatives $D_y^\alpha D_\xi^\alpha$ for decay, and noting
that the first argument is still $x$ and $|x-(x+t(y-x))|\le|x-y|$, we find
\begin{equation*}
  |D_x^\mu D_y^\nu D_\xi^\gamma r_N|\le C\ang{(x,\xi)}^{m-2N}\ang{x-y}^{w+2N}\ang{\xi}^{-|\gamma|}.
\end{equation*}
Let $2N>m$ and $s=N-m/2$. By \eqref{isotropic:eq:split},
$\ang{(x,\xi)}^{-2s}\le\ang{x}^{-s}\ang{\xi}^{-s}$, so $\ang{x}^sr_N$ is an amplitude
of the traditional calculus of order $-s$ with weight $\ang{x-y}^{w+2N}$
(derivatives of $\ang{x}^s$ are bounded by $C\ang{x}^s$). By
\cref{ch:euclidean}, $\Op(r_N)=\ang{x}^{-s}Q_N$ with $Q_N\in\Psido{-s}(\R^n)$,
so the left symbol of $\Op(r_N)$ lies in $\ang{x}^{-s}S^{-s}(\R^n;\R^n)$: all its
derivatives are bounded by $C\ang{x}^{-s}\ang{\xi}^{-s}$.

Now $\Op(a)$ has a left symbol $\sigma\in\ang{x}^{m_+}S^m(\R^n;\R^n)$ and, by
\eqref{isotropic:eq:taylor-remainder} and uniqueness of left symbols,
$\sigma-\sum_{|\alpha|<N}c_\alpha\in\ang{x}^{-s}S^{-s}$. Choose
$b\in\Gs{m}(\R^{2n})$ with $b\sim\sum_\alpha c_\alpha$
(\cref{isotropic:lem:asymp-Gamma}). Then for every $N$
\begin{equation*}
  \sigma-b=\Bigl(\sigma-\sum_{|\alpha|<N}c_\alpha\Bigr)-\Bigl(b-\sum_{|\alpha|<N}c_\alpha\Bigr)
  \in\ang{x}^{-s}S^{-s}+\Gs{m-2N},
\end{equation*}
and all derivatives of elements of both spaces are bounded by
$C\ang{x}^{-s}\ang{\xi}^{-s}$, with $s=N-m/2$ arbitrarily large. Hence
$\sigma-b\in\cS(\R^{2n})$, so $\sigma\in\Gs{m}(\R^{2n})$ and
$\sigma\sim\sum_\alpha c_\alpha$.
\end{proof}

\subsection{The algebra of isotropic operators}

\begin{theorem}\label{isotropic:thm:algebra}
The spaces $\Psiiso{m}(\R^n)$ form an order-filtered $*$-algebra with residual
algebra $\Psiiso{-\infty}(\R^n)$, the operators with kernels in
$\cS(\R^{2n})$. Precisely, if $A=\Op(a)\in\Psiiso{m}(\R^n)$ and
$B=\Op(b)\in\Psiiso{m'}(\R^n)$ then
\begin{enumerate}[label=\textup{(\roman*)}]
\item the formal adjoint $A^*$, defined by $(Au,v)=(u,A^*v)$ for
  $u,v\in\cS(\R^n)$, lies in $\Psiiso{m}(\R^n)$ and
  $\sigma_L(A^*)\sim\sum_\alpha\frac{i^{|\alpha|}}{\alpha!}D_\xi^\alpha D_x^\alpha\bar a$;
\item $AB\in\Psiiso{m+m'}(\R^n)$ and
  \begin{equation}\label{isotropic:eq:composition}
    \sigma_L(AB)=a\# b\sim\sum_\alpha\frac1{\alpha!}\pa_\xi^\alpha a\,D_x^\alpha b,
    \qquad a\#b-ab\in\Gs{m+m'-2}(\R^{2n});
  \end{equation}
\item the subspaces $\Psiisoph{m}(\R^n)$ form a $*$-subalgebra;
\item the bilinear maps $\Gs{m}\times\Gs{m'}\ni(a,b)\mapsto a\#b\in\Gs{m+m'}$
  and $(a,b)\mapsto a\#b-\sum_{|\alpha|<N}\frac1{\alpha!}\pa_\xi^\alpha a\,D_x^\alpha b\in\Gs{m+m'-2N}$
  are continuous.
\end{enumerate}
\end{theorem}
\begin{proof}
(i) The kernel of $A^*$ is $\overline{K_A(y,x)}=\int e^{i(x-y)\cdot\xi}\bar a(y,\xi)\,\dbar\xi$,
so $A^*=\Op(\bar a(y,\xi))$, and $\bar a(y,\xi)\in\Ampl{m}$ by
\cref{isotropic:lem:ampl-examples}(ii) with $\lambda=1$; apply
\cref{isotropic:prop:reduction}.

(ii) Applying (i) to $B^*=\Op(c)$ we find $B=\Op(c)^*=\Op(b_R(y,\xi))$ with
$b_R=\bar c\in\Gs{m'}$; taking conjugates in the expansion of (i),
$b_R\sim\sum_\gamma\frac{(-i)^{|\gamma|}}{\gamma!}D_\xi^\gamma D_x^\gamma b$. For
$u\in\cS$, $\widehat{Bu}(\xi)=\int e^{-iy\cdot\xi}b_R(y,\xi)u(y)\,dy$, hence
$ABu(x)=\int e^{ix\cdot\xi}a(x,\xi)\widehat{Bu}(\xi)\,\dbar\xi$, i.e.
$AB=\Op(a(x,\xi)b_R(y,\xi))$. By \cref{isotropic:lem:ampl-examples}(i) and
\cref{isotropic:prop:reduction}, $AB\in\Psiiso{m+m'}$ with
$\sigma_L(AB)\sim\sum_\beta\frac{i^{|\beta|}}{\beta!}D_\xi^\beta(a\,D_x^\beta b_R)$.
Inserting the expansion of $b_R$ (with errors in $\Gs{m'-2N}$, which
contribute errors in $\Gs{m+m'-2N}$) we obtain a double series whose terms of
index $(\beta,\gamma)$ lie in $\Gs{m+m'-2|\beta|-2|\gamma|}$. Grouping the
terms with $|\beta|+|\gamma|=j$ we get a finite sum which, by the identity of
formal power series in commuting variables
\begin{equation*}
  e^{iD_y\cdot(D_\eta+D_\xi)}e^{-iD_y\cdot D_\xi}=e^{iD_y\cdot D_\eta}
\end{equation*}
applied to $a(x,\eta)b(y,\xi)$ and evaluated at $y=x$, $\eta=\xi$, equals
$\sum_{|\alpha|=j}\frac{i^{|\alpha|}}{\alpha!}D_\xi^\alpha a\,D_x^\alpha b$.
Since $i^{|\alpha|}D_\xi^\alpha=\pa_\xi^\alpha$ this is
\eqref{isotropic:eq:composition}; the terms with $\alpha\ne0$ lie in
$\Gs{m+m'-2}$.

(iii) All terms in the expansions are polyhomogeneous when $a,b$ are, and
asymptotic sums of polyhomogeneous symbols are polyhomogeneous.

(iv) By the closed graph theorem for Fréchet spaces it suffices to check that
$a\mapsto a\#b$ (for fixed $b$) and $b\mapsto a\#b$ have closed graphs
(separately continuous bilinear maps of Fréchet spaces are continuous). If
$a_k\to a$ in $\Gs{m}$ and $a_k\#b\to c$ in $\Gs{m+m'}$, then for
$u,v\in\cS(\R^n)$, $\pair{\Op(a_k\#b)u}{v}=\pair{\Op(a_k)\Op(b)u}{v}$; the
left side tends to $\pair{\Op(c)u}{v}$ and the right side to
$\pair{\Op(a\#b)u}{v}$, since $\pair{\Op(a)f}{v}=\int\!\!\int e^{ix\cdot\xi}a(x,\xi)\hat f(\xi)v(x)\,dx\,\dbar\xi$
depends continuously on $a\in\Gs{m}$ for $f,v\in\cS$. Thus $c=a\#b$. The second
map is handled in the same way.
\end{proof}

\begin{remark}\label{isotropic:rem:principal}
The \emph{principal symbol} $\sigma_m(A)$ of $A\in\Psiiso{m}(\R^n)$ is the
class of $\sigma_L(A)$ in $\Gs{m}/\Gs{m-1}$; there is a short exact sequence
$0\to\Psiiso{m-1}\to\Psiiso{m}\to\Gs{m}/\Gs{m-1}\to0$, and
$\sigma_{m+m'}(AB)=\sigma_m(A)\sigma_{m'}(B)$, $\sigma_m(A^*)=\overline{\sigma_m(A)}$.
For $A\in\Psiisoph{m}$ we identify $\sigma_m(A)$ with the principal symbol of
$\sigma_L(A)$, a function in $\CI(\Sphere^{2n-1})$. A feature of the isotropic
calculus is that the expansions above proceed in steps of two: $AB-\Op(ab)\in\Psiiso{m+m'-2}$
and, by \eqref{isotropic:eq:composition},
$\sigma_L([A,B])-(-i)\{a,b\}\in\Gs{m+m'-4}$, where
$\{a,b\}=\sum_j(\pa_{\xi_j}a\,\pa_{x_j}b-\pa_{x_j}a\,\pa_{\xi_j}b)$. Under the
identification of \cref{isotropic:prop:rc}, the theorem shows that the
Laurent functions $\mathcal L(\Hsph{2n})$ carry a non-commutative product, for
which $\CI(\Hsph{2n})$ is a subalgebra.
\end{remark}

\begin{remark}\label{isotropic:rem:vector}
If $E$ and $F$ are finite-dimensional vector spaces, $\Psiiso{m}(\R^n;E,F)$
denotes the operators from $\cS(\R^n;E)$ to $\cS(\R^n;F)$ given, after a choice
of bases, by matrices of elements of $\Psiiso{m}(\R^n)$; equivalently their
left symbols are in $\Gs{m}(\R^{2n};\Hom(E,F))$. All results of this chapter
extend to this setting with the obvious changes, and we use them freely; in
\cref{isotropic:sec:L2} we work with $\C^N$-valued operators explicitly.
\end{remark}

\subsection{Invariance under the Fourier transform}

A property which distinguishes the isotropic calculus sharply from the
traditional one is its invariance under the Fourier transform.

\begin{proposition}\label{isotropic:prop:fourier}
If $A\in\Psiiso{m}(\R^n)$ then $\hat A=\cF A\cF^{-1}\in\Psiiso{m}(\R^n)$, i.e.\
$\widehat{Au}=\hat A\hat u$; if $A=\Op(b(y,\xi))$ is written in right-reduced
form, $b\in\Gs{m}(\R^{2n})$, then $\sigma_L(\hat A)(\xi,z)=b(-z,\xi)$. The same holds
for $\Psiisoph{m}(\R^n)$.
\end{proposition}
\begin{proof}
The right-reduced form exists by the proof of
\cref{isotropic:thm:algebra}(ii). Suppose first $b\in\cS(\R^{2n})$. Then
$\widehat{Au}(\xi)=\int e^{-iy\cdot\xi}b(y,\xi)u(y)\,dy$, and substituting
$u(y)=\int e^{iy\cdot\eta}\hat u(\eta)\,\dbar\eta$ and then $y=-z$,
\begin{equation*}
  \widehat{Au}(\xi)=\int\!\!\int e^{i(\xi-\eta)\cdot z}b(-z,\xi)\hat u(\eta)\,d\eta\,\dbar z ,
\end{equation*}
which is $\Op(c)\hat u$ with $c(\xi,z)=b(-z,\xi)$, $\xi$ now playing the role of
the base variable and $z$ that of the dual variable. Clearly $c\in\Gs{m}$ if
$b\in\Gs{m}$, by \cref{isotropic:lem:Gamma-basic}(iii). For general
$b\in\Gs{m}$ choose $b_k\in\cS$ with $b_k\to b$ in $\Gs{m+1}$
(\cref{isotropic:lem:Gamma-basic}(iv)); for $u,v\in\cS$ both
$\pair{\cF\Op(b_k(y,\xi))\cF^{-1}u}{v}$ and $\pair{\Op(c_k)u}{v}$ depend
continuously on $b_k\in\Gs{m+1}$, so the identity persists in the limit.
\end{proof}

%======================================================================
\section{Scattering operators}\label{isotropic:sec:scattering}

\begin{definition}\label{isotropic:def:Psisc}
For $m,l\in\R$ set $\Psisc{m,l}(\R^n)=\{\Op(a);\ a\in S^{m,l}(\R^n;\R^n)\}$.
\end{definition}

Thus $\Psisc{m,l}(\R^n)\subset\ang{x}^l\Psido{m}(\R^n)$, and by
\cref{isotropic:lem:Gamma-in-S}, $\Psiiso{m}\subset\Psisc{m,m}$ if $m\ge0$ and
$\Psiiso{m}\subset\Psisc{m-t,t}$ for $m\le t\le0$. Since
$\bigcap_{m,l}S^{m,l}(\R^n;\R^n)=\cS(\R^{2n})$,
\begin{equation}\label{isotropic:eq:sc-residual}
  \bigcap_{m,l}\Psisc{m,l}(\R^n)=\Psiiso{-\infty}(\R^n).
\end{equation}
Typical elements are the differential operators
$\sum_{|\alpha|\le m}c_\alpha(x)D^\alpha$ with
$|\pa^\beta c_\alpha|\le C_\beta\ang{x}^{l-|\beta|}$; for instance the
Laplacian $\Delta=\sum_jD_j^2$ is in $\Psisc{2,0}(\R^n)$, and $\Delta+1$ is
elliptic in the scattering sense below, although it is not an isotropic
operator.

For scattering operators one uses the amplitudes $\Ampl{m,l}_{\mathrm{sc}}$:
smooth functions $a(x,y,\xi)$ such that for some $w\ge0$
\begin{equation*}
  |D_x^\alpha D_y^\beta D_\xi^\gamma a|\le C\ang{x}^{l-|\alpha'|-|\beta'|}\ang{x-y}^{w+|\alpha'|+|\beta'|}\ang{\xi}^{m-|\gamma|}
  \qquad\forall\,\alpha'\le\alpha,\ \beta'\le\beta .
\end{equation*}
As in \cref{isotropic:lem:ampl-examples}, $a(x,\xi)b(y,\xi)\in\Ampl{m+m',l+l'}_{\mathrm{sc}}$
if $a\in S^{m,l}$, $b\in S^{m',l'}$ (use \eqref{isotropic:eq:peetre} for
$\ang{y}^{l'-|\beta'|}$), and $a(x+\lambda(y-x),\xi)\in\Ampl{m,l}_{\mathrm{sc}}$.

\begin{theorem}\label{isotropic:thm:sc-algebra}
The spaces $\Psisc{m,l}(\R^n)$ form a bi-filtered $*$-algebra,
$\Psisc{m,l}\circ\Psisc{m',l'}\subset\Psisc{m+m',l+l'}$, with residual space
\eqref{isotropic:eq:sc-residual}. The left symbols of adjoints and products are
given by the formul\ae\ of \cref{isotropic:thm:algebra}, the expansions now being
asymptotic in the sense of \cref{isotropic:lem:asymp-sc}; in particular
$AB-\Op(ab)\in\Psisc{m+m'-1,l+l'-1}(\R^n)$.
\end{theorem}
\begin{proof}
The proof of \cref{isotropic:prop:reduction} applies with the following
changes. For $a\in\Ampl{m,l}_{\mathrm{sc}}$ the terms
$c_\alpha=\frac{i^{|\alpha|}}{\alpha!}(D_\xi^\alpha D_y^\alpha a)(x,x,\xi)$ lie in
$S^{m-|\alpha|,l-|\alpha|}$, and the remainder in
\eqref{isotropic:eq:taylor-remainder} satisfies
$|D_x^\mu D_y^\nu D_\xi^\gamma r_N|\le C\ang{x}^{l-N}\ang{x-y}^{w+N}\ang{\xi}^{m-N-|\gamma|}$.
Hence $\Op(r_N)\in\ang{x}^{l-N}\Psido{m-N}(\R^n)$ and the left symbol $\sigma$
of $\Op(a)$ satisfies $\sigma-\sum_{|\alpha|<N}c_\alpha\in\ang{x}^{l-N}S^{m-N}(\R^n;\R^n)$.
If $b\in S^{m,l}$ with $b\sim\sum c_\alpha$, then
$\sigma-b\in\bigcap_N\ang{x}^{l-N}S^{m-N}(\R^n;\R^n)=\cS(\R^{2n})$. Adjoints and
products are then treated exactly as in \cref{isotropic:thm:algebra}.
\end{proof}

\begin{remark}\label{isotropic:rem:sc-ph}
Polyhomogeneous scattering symbols are defined using the radial
compactification in each factor:
$S^{m,l}_{\mathrm{ph}}(\R^p;\R^n)=\ang{z}^l\ang{\xi}^m(\rc\times\rc)^*\CI(\Hsph{p}\times\Hsph{n})$.
By \cref{isotropic:prop:rc} applied in each factor these are scattering
symbols, and they are preserved by $D_z$, $D_\xi$ (which lower the respective
orders) and by products. If
$\Psisc{m,l}_{\mathrm{ph}}(\R^n)=\Op(S^{m,l}_{\mathrm{ph}})$ then
$\Psisc{m,l}_{\mathrm{ph}}\circ\Psisc{m',l'}_{\mathrm{ph}}\subset\Psisc{m+m',l+l'}_{\mathrm{ph}}$:
the terms of the composition formula are polyhomogeneous, an asymptotic sum
$\sum_j c_j$ with $c_j\in S^{m-j,l-j}_{\mathrm{ph}}$ can be realized by a
polyhomogeneous symbol using Borel's lemma for the series
$\sum_j(\rho_1\rho_2)^jG_j$ on the manifold with corners $\Hsph{p}\times\Hsph{n}$
($\rho_1,\rho_2$ the two boundary defining functions), and $\cS(\R^{p+n})$ is
contained in every $S^{m,l}_{\mathrm{ph}}$. Repeating the computation in the
proof of \cref{isotropic:prop:fourier} shows that
$\cF\,\Psisc{m,l}(\R^n)\,\cF^{-1}=\Psisc{l,m}(\R^n)$: the Fourier transform
interchanges the two orders.
\end{remark}

%======================================================================
\section{The residual isotropic algebra}\label{isotropic:sec:residual}

The residual algebra $\Psiiso{-\infty}(\R^n)$, a subalgebra of the residual
algebra $\Psido{-\infty}(\R^n)$ of the traditional calculus, has two important
properties not shared by the latter: its elements are compact on $L^2(\R^n)$,
and it is a `corner' in the bounded operators. We use the strong (dual)
topology on $\cSp(\R^n)$, in which a set is bounded if and only if it is
weakly bounded (Banach--Steinhaus, $\cS$ being Fréchet), and convergence means
uniform convergence on bounded subsets of $\cS$.

\begin{lemma}\label{isotropic:lem:partial-pairing}
If $u\in\cSp(\R^q)$ and $F\in\cS(\R^p\times\R^q)$ then
$f(z)=\pair{u}{F(z,\cdot)}$ lies in $\cS(\R^p)$, with
$\pa_z^\alpha f=\pair{u}{\pa_z^\alpha F(z,\cdot)}$, and $(u,F)\mapsto f$ is continuous
(hypocontinuous) and bounded on products of bounded sets.
\end{lemma}
\begin{proof}
The map $z\mapsto F(z,\cdot)\in\cS(\R^q)$ is smooth, by Taylor's formula with
remainder and the uniform Schwartz bounds on the derivatives of $F$. The
distribution $u$ has finite order:
$|\pair{u}{\phi}|\le C\sum_{|\beta|+|\gamma|\le k}\sup|y^\beta D^\gamma\phi|$. Hence
$|z^\delta\pa_z^\alpha f(z)|\le C\sum_{|\beta|+|\gamma|\le k}\sup_y|z^\delta y^\beta D_y^\gamma\pa_z^\alpha F(z,y)|$,
which is bounded.
\end{proof}

\begin{proposition}\label{isotropic:prop:S'-S}
A continuous linear operator $A:\cS(\R^n)\to\cSp(\R^n)$ lies in
$\Psiiso{-\infty}(\R^n)$ if and only if it extends to a continuous linear map
\begin{equation*}
  A:\cSp(\R^n)\to\cS(\R^n).
\end{equation*}
In particular each $A\in\Psiiso{-\infty}(\R^n)$ maps bounded subsets of
$L^2(\R^n)$ into bounded subsets of $\cS(\R^n)$.
\end{proposition}
\begin{proof}
If $K_A\in\cS(\R^{2n})$ then $Au(x)=\pair{u}{K_A(x,\cdot)}$ for $u\in\cSp$ by
\cref{isotropic:lem:partial-pairing}, and
$\sup_x|x^\gamma D^\alpha Au(x)|\le\sup_{\phi\in B}|\pair{u}{\phi}|$ where
$B=\{x^\gamma D_x^\alpha K_A(x,\cdot);\ x\in\R^n\}$ is bounded in $\cS$, which is
continuity for the strong topology.

Conversely suppose $A:\cSp\to\cS$ is continuous. The map
$\R^n\ni y\mapsto\delta_y\in\cSp$ is smooth, with $\pa_{y_j}\delta_y=-\pa_j\delta_y$,
since difference quotients of $\phi\in\cS$ converge uniformly for $\phi$ in
bounded sets. For each $N$ the set
$\{\ang{y}^N\pa^\beta\delta_y;\ y\in\R^n,\ |\beta|\le N\}$ is weakly, hence
strongly, bounded in $\cSp$, because
$|\pair{\ang{y}^N\pa^\beta\delta_y}{\phi}|=\ang{y}^N|\pa^\beta\phi(y)|$ is bounded for
each $\phi\in\cS$. Its image under $A$ is bounded in $\cS$. Thus
$K(x,y)=(A\delta_y)(x)$ is smooth and $\ang{y}^N|x^\gamma D_x^\alpha D_y^\beta K(x,y)|$
is bounded for all $N,\gamma,\alpha$ and $|\beta|\le N$, i.e.\
$K\in\cS(\R^{2n})$. Finally, for $\phi\in\cS$ the Riemann sums of
$\int\phi(y)\delta_y\,dy$ converge to $\phi$ in $\cSp$, so
$A\phi=\int\phi(y)A\delta_y\,dy=\int K(\cdot,y)\phi(y)\,dy$ and $K$ is the kernel of $A$.
\end{proof}

\begin{proposition}\label{isotropic:prop:compact}
Each $A\in\Psiiso{-\infty}(\R^n)$ is a Hilbert--Schmidt, hence compact, operator
on $L^2(\R^n)$, and $\norm{A}_{\Lin(L^2)}\le\norm{K_A}_{L^2(\R^{2n})}$; in particular
the inclusion $\Psiiso{-\infty}(\R^n)\to\Lin(L^2(\R^n))$ is continuous.
\end{proposition}
\begin{proof}
$K_A\in\cS(\R^{2n})\subset L^2(\R^{2n})$. (Alternatively: by
\cref{isotropic:prop:S'-S}, $A$ maps the unit ball of $L^2$ into a bounded set
in $\cS$, and bounded subsets of $\cS$ are precompact in $L^2$ by the
Arzelà--Ascoli theorem applied to $\ang{x}^{n+1}u$.)
\end{proof}

\begin{lemma}\label{isotropic:lem:corner}
If $A_1,A_2\in\Psiiso{-\infty}(\R^n)$ and $B$ is a bounded operator on
$L^2(\R^n)$, or more generally any continuous linear operator
$B:\cS(\R^n)\to\cSp(\R^n)$, then $A_1BA_2\in\Psiiso{-\infty}(\R^n)$. Moreover for
each $k$ there are $k'$ and $C$ such that, for bounded $B$,
\begin{equation}\label{isotropic:eq:corner-est}
  p_k(K_{A_1BA_2})\le C\,p_{k'}(K_{A_1})\,p_{k'}(K_{A_2})\,\norm{B}_{\Lin(L^2)},
\end{equation}
where $p_k(K)=\sum_{|\alpha|+|\beta|+|\gamma|+|\delta|\le k}\sup|x^\alpha y^\beta D_x^\gamma D_y^\delta K(x,y)|$.
\end{lemma}
\begin{proof}
For bounded $B$ the kernel of $C=A_1BA_2$ is
\begin{equation*}
  K_C(x,y)=\bigl(B[K_{A_2}(\cdot,y)],\overline{K_{A_1}(x,\cdot)}\bigr)_{L^2},
\end{equation*}
and differentiating under the inner product (the maps
$y\mapsto K_{A_2}(\cdot,y)$, $x\mapsto K_{A_1}(x,\cdot)$ are smooth into $L^2$)
\begin{equation*}
  |x^\alpha y^\beta D_x^\gamma D_y^\delta K_C(x,y)|\le\norm{B}\,
  \sup_y\norm{y^\beta D_y^\delta K_{A_2}(\cdot,y)}_{L^2}\,\sup_x\norm{x^\alpha D_x^\gamma K_{A_1}(x,\cdot)}_{L^2},
\end{equation*}
and $\norm{f(x,\cdot)}_{L^2}\le C\sup_{x'}\ang{x'}^{n}|f(x,x')|$. For general $B$
with kernel $K_B\in\cSp(\R^{2n})$ (Schwartz kernel theorem),
$K_C(x,y)=\pair{K_B}{K_{A_1}(x,\cdot)\otimes K_{A_2}(\cdot,y)}$, the pairing being in
the variables $(x',y')$ of $K_{A_1}(x,x')K_{A_2}(y',y)\in\cS(\R^{4n})$; this is in
$\cS(\R^{2n})$ by \cref{isotropic:lem:partial-pairing}.
\end{proof}

Note that $\Psiiso{-\infty}(\R^n)$ is an ideal in $\Psiiso{\infty}(\R^n)$, but
not in $\Lin(L^2(\R^n))$.

\subsection{Finite rank operators and perturbations of the identity}

If $B$ is a bounded operator on $L^2(\R^n)$ of finite rank, let
$f_1,\dots,f_N$ be an orthonormal basis of its range; then
$u\mapsto(Bu,f_j)$ are continuous functionals, given by $g_j\in L^2$, and
$K_B=\sum_jf_j(x)\overline{g_j(y)}$. If $B\in\Psiiso{-\infty}(\R^n)$ the range
of $B$ and of $B^*$ lie in $\cS$, so $f_j,g_j\in\cS(\R^n)$; conversely such
kernels are in $\cS(\R^{2n})$. The finite rank elements of
$\Psiiso{-\infty}(\R^n)$ form an ideal in $\Psiiso{\infty}(\R^n)$.

\begin{lemma}\label{isotropic:lem:index-const}
If $K$ is a compact operator on a Hilbert space $\mathcal H$ then for each $t\in\C$,
$\Id+tK$ has finite-dimensional null space and closed range of finite
codimension, and $\ind(\Id+tK)=\dim\Nul(\Id+tK)-\dim\Nul(\Id+\bar tK^*)=0$.
\end{lemma}
\begin{proof}
Let $T=\Id+K$ (replace $K$ by $tK$). On $N_0=\Nul(T)$ we have $u=-Ku$, so the
unit ball of $N_0$ is contained in $K(\text{unit ball})$, which is precompact;
hence $\dim N_0<\infty$. Let $R_0=N_0^\perp$ and $R_1=\Ran(T)$. If
$f_k=Tu_k\to f$ with $u_k\in R_0$, then $(u_k)$ is bounded: otherwise along a
subsequence $\norm{u_k}\to\infty$, and $v_k=u_k/\norm{u_k}$ satisfies
$v_k=-Kv_k+f_k/\norm{u_k}$; by compactness a subsequence converges to some
$v\in R_0$ with $\norm{v}=1$ and $Tv=0$, a contradiction. Then
$u_k=-Ku_k+f_k$ has a convergent subsequence, with limit $u$ satisfying
$Tu=f$; so $R_1$ is closed, and $N_1=R_1^\perp=\Nul(T^*)$ is finite-dimensional
by the first argument applied to $K^*$.

To see that the index is locally constant in $t$, fix $t_0$, let $T=\Id+t_0K$
and write $\Id+tK=T+(t-t_0)K$ as a $2\times2$ matrix with respect to the
decompositions $\mathcal H=N_0\oplus R_0$ (domain) and $\mathcal H=N_1\oplus R_1$
(range):
\begin{equation*}
  \Id+tK=\begin{pmatrix}(t-t_0)K_{00}&(t-t_0)K_{01}\\(t-t_0)K_{10}&D(t)\end{pmatrix},
  \qquad D(t)=\tilde T+(t-t_0)K_{11},
\end{equation*}
where $\tilde T:R_0\to R_1$ is the restriction of $T$, invertible by the closed
graph theorem. For $|t-t_0|$ small $D(t)$ is invertible and
\begin{equation*}
  \begin{pmatrix}P&Q\\R&D\end{pmatrix}=\begin{pmatrix}\Id&QD^{-1}\\0&\Id\end{pmatrix}
  \begin{pmatrix}P-QD^{-1}R&0\\0&D\end{pmatrix}\begin{pmatrix}\Id&0\\D^{-1}R&\Id\end{pmatrix}.
\end{equation*}
The outer factors are invertible, so null space and cokernel of $\Id+tK$ are
isomorphic to those of the linear map $G(t)=P-QD^{-1}R:N_0\to N_1$ between
finite-dimensional spaces, whose index is $\dim N_0-\dim N_1=\ind(\Id+t_0K)$.
Thus $t\mapsto\ind(\Id+tK)$ is locally constant on $\C$, hence constant, and it
vanishes at $t=0$.
\end{proof}

\begin{proposition}\label{isotropic:prop:Id+A}
Let $A\in\Psiiso{-\infty}(\R^n)$. Then $\Id+A$ is Fredholm of index zero on
$L^2(\R^n)$: its null space $N_0$ is finite-dimensional and contained in
$\cS(\R^n)$, and its range is closed and is the orthocomplement of the null
space $N_1\subset\cS(\R^n)$ of $\Id+A^*$. If $\Pi_0,\Pi_1$ denote the orthogonal
projections onto $N_0,N_1$, then $\Pi_0,\Pi_1\in\Psiiso{-\infty}(\R^n)$ and there
is $B\in\Psiiso{-\infty}(\R^n)$ such that
\begin{equation}\label{isotropic:eq:gen-inverse-residual}
  (\Id+A)(\Id+B)=\Id-\Pi_1,\qquad(\Id+B)(\Id+A)=\Id-\Pi_0 .
\end{equation}
Moreover there is a finite rank $A'\in\Psiiso{-\infty}(\R^n)$, of rank
$\dim N_0$, such that $\Id+A+sA'$ is invertible on $L^2(\R^n)$ for all $s\ne0$.
\end{proposition}
\begin{proof}
By \cref{isotropic:prop:compact} and \cref{isotropic:lem:index-const}, $\Id+A$
is Fredholm of index zero. If $u\in N_0$ then $u=-Au\in\cS$, and similarly
$N_1\subset\cS$, so $\Pi_0=\sum_jf_j\otimes\bar f_j$ and $\Pi_1=\sum_jg_j\otimes\bar g_j$,
for orthonormal bases $f_j\in N_0$, $g_j\in N_1$, are in $\Psiiso{-\infty}$. The
restriction of $\Id+A$ is a bijection $N_0^\perp\to N_1^\perp=\Ran(\Id+A)$,
with bounded inverse; extending this inverse by $0$ on $N_1$ gives a bounded
operator $\Id+B$ satisfying \eqref{isotropic:eq:gen-inverse-residual}. From the
first identity $B=-A-AB-\Pi_1$, and from the second
$AB=-A^2-ABA-A\Pi_0$; hence
\begin{equation*}
  B=-A+A^2+ABA+A\Pi_0-\Pi_1\in\Psiiso{-\infty}(\R^n)
\end{equation*}
by \cref{isotropic:lem:corner}. Finally, since $\dim N_0=\dim N_1=N$, the operator
$A'=\sum_{j=1}^Ng_j\otimes\bar f_j$ maps $N_0$ isomorphically onto $N_1$ and
vanishes on $N_0^\perp$; with respect to $L^2=N_0\oplus N_0^\perp\to N_1\oplus N_1^\perp$,
$\Id+A+sA'$ is the direct sum of $sA'|_{N_0}$ and $(\Id+A)|_{N_0^\perp}$, invertible
for $s\ne0$.
\end{proof}

\begin{remark}\label{isotropic:rem:index-trace}
Once the trace on $\Psiiso{-\infty}(\R^n)$ has been introduced
(\cref{ch:determinants}), the vanishing of the index also follows from
\eqref{isotropic:eq:gen-inverse-residual}:
$\ind(\Id+A)=\Tr\Pi_0-\Tr\Pi_1=\Tr([A,B])=0$. The argument above is
independent of the trace.
\end{remark}

\begin{corollary}\label{isotropic:cor:invertible}
If $A\in\Psiiso{-\infty}(\R^n)$ and $\Id+A$ is injective, or surjective, on
$L^2(\R^n)$, then $\Id+A$ is invertible and its inverse lies in
$\Id+\Psiiso{-\infty}(\R^n)$.
\end{corollary}
\begin{proof}
By index zero both $N_0$ and $N_1$ vanish, and $(\Id+A)^{-1}=\Id+B$.
\end{proof}

\begin{corollary}\label{isotropic:cor:Id+A-S}
For $A\in\Psiiso{-\infty}(\R^n)$, $\Id+A$ is Fredholm of index zero as an
operator on $\cS(\R^n)$ and on $\cSp(\R^n)$: in both cases its null space is
$N_0$ and its range is $\{f;\ \Pi_1f=0\}$, a closed subspace of codimension
$\dim N_1$.
\end{corollary}
\begin{proof}
The identities \eqref{isotropic:eq:gen-inverse-residual} hold on $\cS$ and, by
continuity and density, on $\cSp$. If $u\in\cSp$ and $(\Id+A)u=0$ then
$u=-Au\in\cS$, so $u\in N_0$. The range contains $\Ran(\Id-\Pi_1)=\Nul(\Pi_1)$
and is contained in it, since $((\Id+A)u,v)=(u,(\Id+A^*)v)=0$ for $v\in N_1$.
\end{proof}

%======================================================================
\section{Exponential and logarithm}\label{isotropic:sec:exp}

\begin{proposition}\label{isotropic:prop:exp-log}
For $A\in\Psiiso{-\infty}(\R^n)$ the series
\begin{equation*}
  \exp(A)=\sum_{j\ge0}\frac1{j!}A^j
\end{equation*}
converges in $\Id+\Psiiso{-\infty}(\R^n)$, and $\exp:\Psiiso{-\infty}\to\Id+\Psiiso{-\infty}$
is continuous and entire (holomorphic along every complex line, locally
uniformly). Each $\exp(A)$ is invertible with inverse $\exp(-A)$. If
$\norm{A}_{\Lin(L^2)}<1$ the series
\begin{equation*}
  \log(\Id+A)=\sum_{j\ge1}\frac{(-1)^{j+1}}{j}A^j
\end{equation*}
converges in $\Psiiso{-\infty}(\R^n)$, $\exp(\log(\Id+A))=\Id+A$, and
$\log(\exp(A))=A$ if $\norm{A}_{\Lin(L^2)}<\log2$. In particular the range of
$\exp$ contains the neighbourhood $\{\Id+A;\ \norm{A}_{\Lin(L^2)}<1\}$ of $\Id$
in $\Id+\Psiiso{-\infty}(\R^n)$.
\end{proposition}
\begin{proof}
For $j\ge2$, $A^j=A\,A^{j-2}A$ (with $A^0=\Id$), so by
\eqref{isotropic:eq:corner-est}
$p_k(K_{A^j})\le C\,p_{k'}(K_A)^2\norm{A}^{j-2}_{\Lin(L^2)}$. Thus both series
converge absolutely in every seminorm, uniformly for $A$ in sets on which
$p_{k'}(K_A)$ and $\norm{A}$ are bounded (respectively $\norm{A}\le r<1$), which by
\cref{isotropic:prop:compact} include neighbourhoods of each point. The
partial sums are continuous polynomial maps, so the limits are continuous and
holomorphic along complex lines. The identities $\exp(A)\exp(-A)=\Id$,
$\exp(\log(\Id+A))=\Id+A$ and $\log(\exp A)=A$ (where
$\norm{\exp A-\Id}\le e^{\norm{A}}-1<1$) hold for the norm-convergent series in
the Banach algebra $\Lin(L^2)$, and convergence in $\Psiiso{-\infty}$ implies
convergence in $\Lin(L^2)$.
\end{proof}

%======================================================================
\section{Elliptic operators}\label{isotropic:sec:elliptic}

\begin{definition}\label{isotropic:def:elliptic}
An operator $A\in\Psiiso{m}(\R^n;\C^N)$ with left symbol $a$ is \emph{elliptic}
if there are $c,R>0$ such that $a(w)$ is invertible and
$|a(w)^{-1}|\le c^{-1}\ang{w}^{-m}$ for $|w|\ge R$ (for $N=1$:
$|a(w)|\ge c\ang{w}^m$). For $A\in\Psiisoph{m}$ this holds if and only if the
principal symbol $\sigma_m(A)$ is invertible at every point of
$\Sphere^{2n-1}$. Similarly $A\in\Psisc{m,l}(\R^n)$ is elliptic if
$|a(x,\xi)|\ge c\ang{x}^l\ang{\xi}^m$ for $|x|+|\xi|\ge R$.
\end{definition}

\begin{theorem}\label{isotropic:thm:parametrix}
If $A\in\Psiiso{m}(\R^n)$ is elliptic there is $B\in\Psiiso{-m}(\R^n)$, unique
modulo $\Psiiso{-\infty}(\R^n)$, with
\begin{equation}\label{isotropic:eq:parametrix}
  AB-\Id,\ BA-\Id\in\Psiiso{-\infty}(\R^n);
\end{equation}
if $A\in\Psiisoph{m}$ then $B\in\Psiisoph{-m}$. Consequently, if $u\in\cSp(\R^n)$
and $Au\in\cS(\R^n)$ then $u\in\cS(\R^n)$. The same holds for elliptic
$A\in\Psisc{m,l}(\R^n)$, with $B\in\Psisc{-m,-l}(\R^n)$.
\end{theorem}
\begin{proof}
Let $\chi\in\CIc(\R^{2n})$ equal $1$ on $|w|\le R$. Choose $F\in\CI(\C)$ with
$F(z)=1/z$ for $|z|\ge c$; then $F(\ang{w}^{-m}a)\in\Gs{0}$ by
\cref{isotropic:lem:Gamma-basic}(ii), and
$b_0=(1-\chi)\ang{w}^{-m}F(\ang{w}^{-m}a)\in\Gs{-m}$ satisfies $ab_0=1-\chi$.
(For matrix symbols use a smooth function equal to the inverse on the closed
set of matrices $M$ with $|M^{-1}|\le c^{-1}$.) If $a$ is polyhomogeneous so is
$b_0$, since $1/a=a_m^{-1}\sum_k(-r)^k$ asymptotically, with $r=(a-a_m)/a_m\in\Gph{-1}$
for large $|w|$. By \eqref{isotropic:eq:composition}, $B_0=\Op(b_0)$ satisfies
$AB_0=\Id-E$ with $\sigma_L(E)=\chi+\Gs{-2}$, so $E\in\Psiiso{-2}$. Choose
$B_R\in\Psiiso{-m}$ with $\sigma_L(B_R)\sim\sum_j\sigma_L(B_0E^j)$, the terms
having orders $-m-2j$. Then for each $N$, $B_R-\sum_{j<N}B_0E^j\in\Psiiso{-m-2N}$
and
\begin{equation*}
  AB_R=\sum_{j<N}(\Id-E)E^j+\Psiiso{-2N}=\Id-E^N+\Psiiso{-2N}=\Id+\Psiiso{-2N}.
\end{equation*}
Similarly $B_0A=\Id-E'$ with $E'\in\Psiiso{-2}$, and $B_L\sim\sum_jE'^jB_0$
satisfies $B_LA-\Id\in\Psiiso{-\infty}$. Then
$B_L-B_R=B_L(\Id-AB_R)+(B_LA-\Id)B_R\in\Psiiso{-\infty}$, so $B=B_R$ satisfies
\eqref{isotropic:eq:parametrix}; uniqueness follows in the same way. If
$Au\in\cS$ then $u=B(Au)-(BA-\Id)u\in\cS$, since residual operators map $\cSp$
into $\cS$ (\cref{isotropic:prop:S'-S}). In the scattering case the same
construction works with $E\in\Psisc{-1,-1}$ and asymptotic sums in the sense
of \cref{isotropic:lem:asymp-sc}, the residual space being
\eqref{isotropic:eq:sc-residual}.
\end{proof}

Combining the symbolic construction with the analysis of
\cref{isotropic:sec:residual} we see that elliptic operators have generalized
inverses with finite rank errors.

\begin{proposition}\label{isotropic:prop:gen-inverse}
Let $A\in\Psiiso{m}(\R^n)$ be elliptic. Then $N_0=\Nul(A)$ and $N_1=\Nul(A^*)$
(on $\cSp(\R^n)$) are finite-dimensional subspaces of $\cS(\R^n)$, and if
$\Pi_0,\Pi_1$ are the $L^2$-orthogonal projections onto them, there is
$G\in\Psiiso{-m}(\R^n)$ such that
\begin{equation}\label{isotropic:eq:gen-inverse}
  AG=\Id-\Pi_1,\qquad GA=\Id-\Pi_0 .
\end{equation}
Hence $A$ is Fredholm on $\cS(\R^n)$ and on $\cSp(\R^n)$, with range
$\{f;\ \Pi_1f=0\}$ and index $\dim N_0-\dim N_1$. If $A$ is polyhomogeneous
then so is $G$. The same holds for elliptic $A\in\Psisc{m,l}(\R^n)$, with
$G\in\Psisc{-m,-l}(\R^n)$.
\end{proposition}
\begin{proof}
Let $B$ be a parametrix, $AB=\Id-E_R$, $BA=\Id-E_L$ with
$E_R,E_L\in\Psiiso{-\infty}$. If $Au=0$, $u\in\cSp$, then $u=E_Lu$, so
$N_0\subset\Nul(\Id-E_L)$, which is finite-dimensional and contained in $\cS$
by \cref{isotropic:cor:Id+A-S}; the same applies to $A^*$, which is elliptic.

We show that $\Ran_{\cS}(A)=\{f\in\cS;\ f\perp N_1\}$. Applying
\cref{isotropic:prop:Id+A} to $-E_R$ gives $B'\in\Psiiso{-\infty}$ and a finite
rank orthogonal projection $\Pi$, with range $W\subset\cS$, such that
$AB(\Id+B')=(\Id-E_R)(\Id+B')=\Id-\Pi$. Hence
$V=A(\cS)\supset\{f\in\cS;\ f\perp W\}$. Since $W\subset\cS$ is
finite-dimensional, every $v\in V$ decomposes as $v=(v-P_Wv)+P_Wv$ with
$v-P_Wv\in V$, so $V=(W^\perp\cap\cS)\oplus(V\cap W)$, and with
$W_1=W\ominus(V\cap W)$ we get $V=\{f\in\cS;\ f\perp W_1\}$. On the other hand
$W_1=\{g\in\cS;\ g\perp V\}$ (using the density of $W^\perp\cap\cS$ in
$W^\perp$), and $g\in\cS$ is orthogonal to $A(\cS)$ if and only if $A^*g=0$.
Thus $W_1=N_1$, as claimed.

Next, if $f\in\cSp$ and $\Pi_1f=0$, then $f=AB(\Id+B')f+\Pi f$ with $\Pi f\in\cS$
and $\Pi f\perp N_1$ (as $(Au,v)=(u,A^*v)=0$ for $v\in N_1$), so $\Pi f\in A(\cS)$
and $f\in A(\cSp)$. Conversely $A(\cSp)\perp N_1$. Hence we may define
$G:\cSp\to\cSp$ by $G=0$ on $N_1$ and, for $\Pi_1f=0$, $Gf=u$ where $u$ is
the unique solution of $Au=f$ with $\Pi_0u=0$. Then
\eqref{isotropic:eq:gen-inverse} holds on $\cSp$, and $G$ maps $\cS$ to $\cS$
(elliptic regularity). The restriction $G:\cS\to\cS$ has closed graph, hence
is continuous: if $f_k\to f$ and $Gf_k\to u$ in $\cS$ then $Au=(\Id-\Pi_1)f$
and $\Pi_0u=0$, so $u=Gf$. Multiplying \eqref{isotropic:eq:gen-inverse} by $B$
gives $(\Id-E_L)G=B-B\Pi_1$ and $G(\Id-E_R)=B-\Pi_0B$, so
\begin{equation*}
  G=B-B\Pi_1+E_LG=B-B\Pi_1+E_LB-E_L\Pi_0B+E_LGE_R .
\end{equation*}
Here $E_LGE_R$ is continuous from $\cSp$ to $\cS$, since $E_R:\cSp\to\cS$,
$G:\cS\to\cS$ and $E_L:\cS\to\cS$ are; so it is in $\Psiiso{-\infty}$ by
\cref{isotropic:prop:S'-S}. All other terms except $B$ are residual, so
$G\in B+\Psiiso{-\infty}\subset\Psiiso{-m}$. The statements about $\cS$ follow
in the same way, and the scattering case is identical.
\end{proof}

%======================================================================
\section{\texorpdfstring{$L^2$}{L2} boundedness and compactness}\label{isotropic:sec:L2}

\begin{proposition}\label{isotropic:prop:L2}
Every $A\in\Psiiso{0}(\R^n;\C^N)$ is bounded on $L^2(\R^n;\C^N)$, and every
$A\in\Psiiso{-\eps}(\R^n;\C^N)$, $\eps>0$, is compact. Elements of $\Psiiso{m}$
with $m<-3n/2$ are Hilbert--Schmidt.
\end{proposition}
\begin{proof}
By \cref{isotropic:lem:Gamma-in-S}, $\Psiiso{0}(\R^n)\subset\Psido{0}(\R^n)$,
which is bounded on $L^2$ by \cref{ch:euclidean}. If $a\in\Gs{-M}$, $M>n$, the
kernel of $\Op(a)$ satisfies
$(x-y)^\beta K(x,y)=\int e^{i(x-y)\cdot\xi}(-D_\xi)^\beta a(x,\xi)\,\dbar\xi$, and
$\int\ang{(x,\xi)}^{-M}d\xi=C\ang{x}^{n-M}$ (substitute $\xi=\ang{x}\eta$), so
$|K(x,y)|\le C_N\ang{x}^{n-M}\ang{x-y}^{-N}$ for all $N$; this is in
$L^2(\R^{2n})$ if $M>3n/2$. Now let $A\in\Psiiso{-\eps}$ and $T=A^*A\in\Psiiso{-2\eps}$,
a bounded self-adjoint operator. For $j$ large, $T^{2^j}\in\Psiiso{-2^{j+1}\eps}$
is Hilbert--Schmidt, hence compact. If $S$ is bounded and self-adjoint with $S^2$
compact, then $S$ is compact: for $u_k\rightharpoonup u$ weakly,
$\norm{S(u_k-u)}^2=(S^2(u_k-u),u_k-u)\le\norm{S^2(u_k-u)}\norm{u_k-u}\to0$. By
induction $T$ is compact, and then
$\norm{A(u_k-u)}^2=(T(u_k-u),u_k-u)\to0$, so $A$ is compact.
\end{proof}

We next identify the norm of $A\in\Psiiso{0}$ modulo compact operators. The
lower bound uses \emph{coherent states}. For $w=(x_0,\xi_0)\in\R^{2n}$ let
$(T_wu)(x)=e^{ix\cdot\xi_0}u(x-x_0)$, a unitary operator on $L^2(\R^n)$
preserving $\cS(\R^n)$; a direct change of variables shows
\begin{equation}\label{isotropic:eq:phase-translation}
  T_w^{-1}\Op(a)T_w=\Op(a(\cdot+w)).
\end{equation}
Let $u_0=\pi^{-n/4}e^{-|x|^2/2}$, so $\norm{u_0}=1$.

\begin{lemma}\label{isotropic:lem:coherent}
Let $w_k\in\R^{2n}$ with $|w_k|\to\infty$ and $\phi_k=T_{w_k}u_0$. Then
$\norm{\phi_k}=1$, $\phi_k\rightharpoonup0$ weakly in $L^2(\R^n)$, and for every
$a\in\Gs{0}(\R^{2n};M(N,\C))$
\begin{equation*}
  \sup_{v\in\C^N,|v|\le1}\norm{\Op(a)(\phi_kv)-\phi_k\,a(w_k)v}_{L^2}\to0 .
\end{equation*}
\end{lemma}
\begin{proof}
By \eqref{isotropic:eq:phase-translation} the norm in question is
$\norm{\Op(b_k)(u_0v)}$ with $b_k(w)=a(w+w_k)-a(w_k)$. All derivatives of $b_k$
are bounded uniformly in $k$ (as $\Gs{0}\subset\CIb$), and
$|b_k(w)|\le|w|\sup_{t\in[0,1]}|\nabla a(w_k+tw)|\le C|w|\sup_t\ang{w_k+tw}^{-1}\to0$
for each $w$. Now $\Op(b_k)u_0(x)=\int e^{ix\cdot\xi}b_k(x,\xi)\hat u_0(\xi)\,\dbar\xi$;
integrating by parts,
$x^\beta\Op(b_k)u_0=\int e^{ix\cdot\xi}(-D_\xi)^\beta(b_k\hat u_0)\,\dbar\xi$, so
$|\Op(b_k)u_0(x)|\le C\ang{x}^{-n-1}$ uniformly in $k$, while
$\Op(b_k)u_0(x)\to0$ pointwise by dominated convergence. Dominated
convergence in $L^2$ gives $\norm{\Op(b_k)u_0}\to0$.

For the weak convergence it suffices, by density, to show
$(\phi_k,f)\to0$ for $f\in\cS$, and for this it suffices to show that every
subsequence has a further subsequence along which this holds. Write
$w_k=(x_k,\xi_k)$. If $|x_k|\to\infty$ along a subsequence then
$|(\phi_k,f)|\le\int u_0(x-x_k)|f(x)|\,dx\to0$. Otherwise $x_k\to x_*$ along a
subsequence and $|\xi_k|\to\infty$; then $(\phi_k,f)=\hat g_k(-\xi_k)$ with
$g_k=u_0(\cdot-x_k)\bar f\to g_*$ in $L^1$, and
$|\hat g_k(-\xi_k)|\le\norm{g_k-g_*}_{L^1}+|\hat g_*(-\xi_k)|\to0$ by the
Riemann--Lebesgue lemma.
\end{proof}

\begin{theorem}\label{isotropic:thm:ess-norm}
If $A\in\Psiiso{0}(\R^n;\C^N)$ has left symbol $a$ then
\begin{equation}\label{isotropic:eq:ess-norm}
  \inf_{K\ \mathrm{compact}}\norm{A+K}_{\Lin(L^2)}
  =\inf_{E\in\Psiiso{-\infty}(\R^n;\C^N)}\norm{A+E}_{\Lin(L^2)}
  =\limsup_{|w|\to\infty}\norm{a(w)}.
\end{equation}
If $A\in\Psiisoph{0}$ the right side is $\sup_{p\in\Sphere^{2n-1}}\norm{\sigma_0(A)(p)}$.
\end{theorem}
\begin{proof}
Let $\mu$ be the right side of \eqref{isotropic:eq:ess-norm}. Since residual
operators are compact, the first infimum is at most the second.

\emph{Lower bound.} Let $K$ be compact. Choose $w_k\to\infty$ with
$\norm{a(w_k)}\to\mu$ and unit vectors $v_k$ with $|a(w_k)v_k|=\norm{a(w_k)}$;
passing to a subsequence $v_k\to v$. With $\phi_k$ as in
\cref{isotropic:lem:coherent}, $K(\phi_kv)\to0$ since $K$ is compact and
$\phi_k\rightharpoonup0$, so
\begin{equation*}
  \norm{A+K}\ge\liminf_k\norm{(A+K)(\phi_kv)}\ge\liminf_k|a(w_k)v|=\mu .
\end{equation*}

\emph{Upper bound.} Let $C>\mu^2$ and choose $R$ and $\delta'>0$ such that
$p=C-a^*a\ge2\delta'$ for $|w|\ge R$. With $\chi\in\CIc$ equal to $1$ on
$|w|\le R$, $0\le\chi\le1$ and supported in $|w|\le R+1$, the function
$p'=(1-\chi)p+\chi C$ lies in $\Gs{0}$, satisfies $p'\ge2\delta'$ everywhere and
equals $p$ for $|w|\ge R+1$. Since the square root is smooth on positive
definite matrices, $q_0=(p')^{1/2}\in\Gs{0}$ is hermitian, invertible, with
$q_0^{-1}\in\Gs{0}$ (\cref{isotropic:lem:Gamma-basic}(ii)). Let $Q_0=\Op(q_0)$.
By \cref{isotropic:thm:algebra}, $G_1=Q_0^*Q_0-C+A^*A$ is self-adjoint and lies
in $\Psiiso{-2}$, since its symbol is $p'-p$ (compactly supported) modulo
$\Gs{-2}$. Suppose $Q=Q_0+R_1+\dots+R_{k-1}$ has been found with $R_j\in\Psiiso{-2j}$
and $G_k=Q^*Q-C+A^*A\in\Psiiso{-2k}$, and let $g_k=\sigma_L(G_k)$, so that
$g_k-g_k^*\in\Gs{-2k-2}$. Set $R_k=\Op(r_k)$, $r_k=-\tfrac12q_0^{-1}g_k$. Then
$(Q+R_k)^*(Q+R_k)-C+A^*A=G_k+Q^*R_k+R_k^*Q+R_k^*R_k$, whose symbol is, modulo
$\Gs{-2k-2}$,
\begin{equation*}
  g_k+q_0r_k+r_k^*q_0=g_k-\tfrac12g_k-\tfrac12g_k^*\equiv0 .
\end{equation*}
Taking $Q\in\Psiiso{0}$ with $\sigma_L(Q)\sim q_0+\sum_k r_k$ we obtain
$Q^*Q=C-A^*A+G$ with $G\in\Psiiso{-\infty}$ self-adjoint. Hence for $u\in\cS$
\begin{equation*}
  \norm{Au}^2=C\norm{u}^2-\norm{Qu}^2+(Gu,u)\le C\norm{u}^2+(Gu,u).
\end{equation*}
Given $\delta>0$, approximate the compact operator $G$ within $\delta/2$ by a
finite rank operator $\sum_{j\le J}f_j\otimes\bar g_j$ and then replace the $g_j$
by nearby $g_j'\in\cS$, to get a finite-dimensional $W=\operatorname{span}\{g_j'\}\subset\cS$
with $\norm{Gu}\le\delta\norm{u}$ for $u\perp W$. Let $\Pi_W\in\Psiiso{-\infty}$ be the
orthogonal projection onto $W$. Then for $u\in\cS$, $u'=(\Id-\Pi_W)u\perp W$ and
$\norm{A(\Id-\Pi_W)u}^2\le(C+\delta)\norm{u}^2$. Since $A\Pi_W\in\Psiiso{-\infty}$
and $C>\mu^2$, $\delta>0$ are arbitrary, the second infimum is at most $\mu$.

For polyhomogeneous $A$, $a-a_0\to0$ at infinity, $a_0$ being homogeneous of
degree $0$ outside the unit ball, so $\mu=\sup_{\Sphere^{2n-1}}\norm{a_0}$.
\end{proof}

\begin{corollary}\label{isotropic:cor:compact-iff}
An operator $A\in\Psiiso{0}(\R^n;\C^N)$ is compact on $L^2(\R^n;\C^N)$ if and
only if $\sigma_L(A)(w)\to0$ as $|w|\to\infty$. In particular
$A\in\Psiisoph{0}(\R^n;\C^N)$ is compact if and only if $\sigma_0(A)=0$, that is
$A\in\Psiisoph{-1}(\R^n;\C^N)$.
\end{corollary}
\begin{proof}
If $A$ is compact the left side of \eqref{isotropic:eq:ess-norm} vanishes.
Conversely if the right side vanishes, $A$ is a norm limit of residual, hence
compact, operators.
\end{proof}

\begin{corollary}\label{isotropic:cor:fredholm-L2}
An operator $A\in\Psiiso{0}(\R^n;\C^N)$ is Fredholm on $L^2(\R^n;\C^N)$ if and
only if it is elliptic.
\end{corollary}
\begin{proof}
If $A$ is elliptic, \eqref{isotropic:eq:gen-inverse} with $G\in\Psiiso{0}$
bounded shows that $A$ is Fredholm. Conversely, if $A$ is Fredholm there is a
bounded $B$ with $BA=\Id-P$, $P$ of finite rank. If $A$ is not elliptic there
are $w_k\to\infty$ and unit vectors $v_k$ with $|a(w_k)v_k|\to0$; passing to a
subsequence $v_k\to v$. By \cref{isotropic:lem:coherent}, $\norm{A(\phi_kv)}\to0$,
hence $\phi_kv=BA(\phi_kv)+P(\phi_kv)\to0$, since $P$ is compact and
$\phi_k\rightharpoonup0$. This contradicts $\norm{\phi_kv}=1$.
\end{proof}

In particular ellipticity is necessary in \cref{isotropic:prop:gen-inverse}.

%======================================================================
\section{The harmonic oscillator}\label{isotropic:sec:harmonic}

The \emph{harmonic oscillator} is the differential operator
\begin{equation*}
  H=\sum_{j=1}^n(D_j^2+x_j^2)=\Delta+|x|^2 .
\end{equation*}
It is an elliptic element of $\Psiisoph{2}(\R^n)$ with symbol $|w|^2$. Its
spectral decomposition is computed explicitly using the \emph{creation} and
\emph{annihilation} operators
\begin{equation*}
  C_j=D_j+ix_j,\qquad A_j=D_j-ix_j=C_j^*,\qquad j=1,\dots,n .
\end{equation*}

\begin{lemma}\label{isotropic:lem:HO-identities}
The following identities hold for $j,k=1,\dots,n$:
\begin{gather*}
  [A_j,A_k]=[C_j,C_k]=0,\qquad[A_j,C_k]=2\delta_{jk},\\
  H=\sum_jC_jA_j+n=\sum_jA_jC_j-n,\qquad[C_j,H]=-2C_j,\qquad[A_j,H]=2A_j .
\end{gather*}
Hence $HC_j=C_j(H+2)$ and $HA_j=A_j(H-2)$.
\end{lemma}
\begin{proof}
Since $[D_j,x_k]=-i\delta_{jk}$,
$C_jA_j=D_j^2+x_j^2-i[D_j,x_j]=D_j^2+x_j^2-1$ and $A_jC_j=D_j^2+x_j^2+1$, which
gives the formul\ae\ for $H$ and $[A_j,C_j]=2$; the other commutators of the
$A$'s and $C$'s are clearly zero. Then
$[C_j,H]=\sum_kC_k[C_j,A_k]=-2C_j$ and $[A_j,H]=\sum_k[A_j,C_k]A_k=2A_j$.
\end{proof}

\begin{proposition}\label{isotropic:prop:H-inverse}
For $u\in\cS(\R^n)$,
\begin{equation}\label{isotropic:eq:H-positive}
  (Hu,u)=\sum_j\norm{A_ju}^2+n\norm{u}^2 .
\end{equation}
If $u\in\cSp(\R^n)$ and $(H-\lambda)u=0$ for some $\lambda\in\C$, then
$u\in\cS(\R^n)$, and $u\ne0$ implies $\lambda\in\R$, $\lambda\ge n$. The
operator $H$ is invertible on $\cS(\R^n)$ and on $\cSp(\R^n)$, with inverse
$H^{-1}\in\Psiisoph{-2}(\R^n)$, which is a compact, self-adjoint, injective
operator on $L^2(\R^n)$. More generally $H^p\in\Psiisoph{2p}(\R^n)$ for all
$p\in\Z$.
\end{proposition}
\begin{proof}
\eqref{isotropic:eq:H-positive} follows from $H=\sum C_jA_j+n$ and $C_j^*=A_j$.
Since $H-\lambda$ is elliptic, eigenfunctions in $\cSp$ lie in $\cS$
(\cref{isotropic:thm:parametrix}); then $\lambda\norm{u}^2=(Hu,u)$ is real and
at least $n\norm{u}^2$. Thus $\Nul(H)=\Nul(H^*)=\{0\}$ on $\cSp$, and
\cref{isotropic:prop:gen-inverse} gives $G\in\Psiisoph{-2}$ with $HG=GH=\Id$ on
$\cSp$. By \cref{isotropic:prop:L2}, $G$ is compact on $L^2$; it is injective
since $HG=\Id$, and self-adjoint since $(Gf,g)=(Gf,HGg)=(HGf,Gg)=(f,Gg)$ for
$f,g\in\cS$.
\end{proof}

\begin{proposition}\label{isotropic:prop:HO-spectrum}
The eigenvalues of $H$ (on $\cSp(\R^n)$) are
\begin{equation*}
  \operatorname{spec}(H)=\{n,n+2,n+4,\dots\}=n+2\N_0 .
\end{equation*}
The eigenvalue $n$ has a one-dimensional eigenspace spanned by the
\emph{ground state} $u_0=\pi^{-n/4}\exp(-|x|^2/2)$.
\end{proposition}
\begin{proof}
By \eqref{isotropic:eq:H-positive}, $Hu=nu$ if and only if $A_ju=0$ for all
$j$. Now $A_ju=-i(\pa_j+x_j)u=-ie^{-|x|^2/2}\pa_j(e^{|x|^2/2}u)$, so for a
distribution $u$, $A_ju=0$ for all $j$ means that $e^{|x|^2/2}u$ is constant,
i.e.\ $u\in\C u_0$. Similarly $C_jv=-ie^{|x|^2/2}\pa_j(e^{-|x|^2/2}v)$, so the
only tempered solution of $C_jv=0$ is $v=0$ (as $e^{x_j^2/2}$ is not tempered);
thus each $C_j$ is injective on $\cSp$.

If $Hu=\lambda u$ with $0\ne u\in\cS$, then by \cref{isotropic:lem:HO-identities}
$H(A_ju)=(\lambda-2)A_ju$ and $H(C_ju)=(\lambda+2)C_ju$, with $C_ju\ne0$. So
$\lambda+2\N_0$ consists of eigenvalues. Since all eigenvalues are $\ge n$,
$A^\alpha u=0$ when $\lambda-2|\alpha|<n$. Let $k$ be the largest integer such
that $A^\alpha u\ne0$ for some $|\alpha|=k$, and $v=A^\alpha u$ for such
$\alpha$. Then $A_jv=0$ for all $j$, so $v\in\C u_0$ and $\lambda-2k=n$.
Combined with $Hu_0=nu_0$ and the ladder property, the eigenvalues are
exactly $n+2\N_0$.
\end{proof}

\begin{proposition}\label{isotropic:prop:HO-eigenspaces}
For $k\in\N_0$ let $E_k$ be the eigenspace of $H$ with eigenvalue $n+2k$. Then
the map
\begin{equation*}
  \{\text{polynomials on }\R^n\text{ homogeneous of degree }k\}\ni p\longmapsto p(C)u_0\in E_k
\end{equation*}
is an isomorphism; thus $\{C^\alpha u_0;\ |\alpha|=k\}$ is a basis of $E_k$ and
$\dim E_k=\binom{n+k-1}{k}$.
\end{proposition}
\begin{proof}
The $C_j$ commute, so $p(C)$ is well defined, and by
\cref{isotropic:lem:HO-identities}, $HC^\alpha u_0=C^\alpha(H+2|\alpha|)u_0=(n+2|\alpha|)C^\alpha u_0$.
From $[A_j,C_k]=2\delta_{jk}$ and $A_ju_0=0$ we get
$A_jC^\alpha u_0=2\alpha_jC^{\alpha-e_j}u_0$ (zero if $\alpha_j=0$), hence for
$|\beta|=|\alpha|$
\begin{equation}\label{isotropic:eq:A-C}
  A^\beta C^\alpha u_0=\begin{cases}0,&\beta\ne\alpha,\\2^{|\alpha|}\alpha!\,u_0,&\beta=\alpha .\end{cases}
\end{equation}
Applying $A^\beta$ to $\sum_{|\alpha|=k}c_\alpha C^\alpha u_0=0$ gives
$c_\beta=0$, so the map is injective.

For surjectivity we show by induction on $k$: if $v\in E_k$ and $A^\beta v=0$
for all $|\beta|=k$, then $v=0$. For $k=0$ this is trivial, as $A^0=\Id$.
For $k>0$ and $j$ fixed, $A_jv\in E_{k-1}$ (or is $0$) and
$A^\beta(A_jv)=A^{\beta+e_j}v=0$ for $|\beta|=k-1$, so $A_jv=0$ by induction;
then $v\in E_0=\C u_0$ and $Hv=(n+2k)v$ force $v=0$. Now let $v\in E_k$ be
orthogonal to all $C^\alpha u_0$, $|\alpha|=k$. Then
$(A^\alpha v,u_0)=(v,C^\alpha u_0)=0$, and $A^\alpha v\in E_0=\C u_0$, so
$A^\alpha v=0$ for all $|\alpha|=k$, whence $v=0$. Since $E_k\subset\cS$ is
finite-dimensional (\cref{isotropic:prop:gen-inverse} applied to $H-n-2k$),
$E_k$ is spanned by the $C^\alpha u_0$.
\end{proof}

Eigenfunctions with different eigenvalues are orthogonal, $H$ being symmetric
on $\cS$, and by \eqref{isotropic:eq:A-C},
$(C^\alpha u_0,C^\beta u_0)=(A^\beta C^\alpha u_0,u_0)=\delta_{\alpha\beta}2^{|\alpha|}\alpha!$
for $|\alpha|=|\beta|$, since $\norm{u_0}=1$. Hence the \emph{Hermite functions}
\begin{equation}\label{isotropic:eq:hermite}
  u_\alpha=(2^{|\alpha|}\alpha!)^{-1/2}C^\alpha u_0,\qquad\alpha\in\N_0^n,
\end{equation}
form an orthonormal set in $L^2(\R^n)$, with $Hu_\alpha=(n+2|\alpha|)u_\alpha$.
In one variable, $(\pa-x)f=e^{x^2/2}\pa(e^{-x^2/2}f)$ gives
$C^ke^{-x^2/2}=i^kH_k(x)e^{-x^2/2}$, where
$H_k(x)=(-1)^ke^{x^2}\frac{d^k}{dx^k}e^{-x^2}$ is the $k$-th Hermite polynomial,
and in general $u_\alpha(x)=\prod_ju_{\alpha_j}(x_j)$. One also checks
from \eqref{isotropic:eq:A-C} that
\begin{equation}\label{isotropic:eq:ladder}
  C_ju_\alpha=(2\alpha_j+2)^{1/2}u_{\alpha+e_j},\qquad A_ju_\alpha=(2\alpha_j)^{1/2}u_{\alpha-e_j}.
\end{equation}

\begin{lemma}\label{isotropic:lem:compact-eigen}
If $K$ is a compact self-adjoint operator on a Hilbert space $V\ne\{0\}$ and
$K\ne0$, then $K$ has an eigenvector with non-zero eigenvalue.
\end{lemma}
\begin{proof}
Let $\tau=\norm{K}>0$ and choose unit vectors $\phi_j$ with
$\norm{K\phi_j}\to\tau$. Passing to a subsequence, $\phi_j\rightharpoonup\phi$
weakly with $\norm{\phi}\le1$, and $K\phi_j\to K\phi$ by compactness. Thus
$\norm{K\phi}=\tau\le\tau\norm{\phi}$, so $\norm{\phi}=1$ and $\phi$ maximizes
$\norm{K\psi}^2=(K^2\psi,\psi)$ on the unit sphere. For $\phi'\perp\phi$ the
function $t\mapsto\norm{K(\phi+t\phi')}^2/\norm{\phi+t\phi'}^2$, $t\in\R$, has a
maximum at $t=0$, so $\operatorname{Re}(K^2\phi,\phi')=0$; replacing $\phi'$ by
$i\phi'$, $(K^2\phi,\phi')=0$. Hence $K^2\phi=\tau^2\phi$, i.e.\
$(K-\tau)(K+\tau)\phi=0$. Either $(K+\tau)\phi\ne0$ is an eigenvector with
eigenvalue $\tau$, or $\phi$ is one with eigenvalue $-\tau$.
\end{proof}

\begin{theorem}\label{isotropic:thm:hermite-basis}
The Hermite functions $u_\alpha$, $\alpha\in\N_0^n$, form an orthonormal basis
of $L^2(\R^n)$; that is, for $u\in L^2(\R^n)$
\begin{equation}\label{isotropic:eq:hermite-expansion}
  u=\sum_\alpha c_\alpha u_\alpha,\qquad c_\alpha=(u,u_\alpha),
\end{equation}
with convergence in $L^2$.
\end{theorem}
\begin{proof}
Let $W$ be the closed span of the $u_\alpha$ and $V=W^\perp$. The operator
$G=H^{-1}$ of \cref{isotropic:prop:H-inverse} satisfies
$Gu_\alpha=(n+2|\alpha|)^{-1}u_\alpha$, so $G(W)\subset W$ and, $G$ being
self-adjoint, $G(V)\subset V$. If $V\ne\{0\}$, the restriction of $G$ to $V$ is
compact, self-adjoint and injective, so by \cref{isotropic:lem:compact-eigen}
there is $0\ne f\in V$ with $Gf=\mu f$, $\mu\ne0$. Then $Hf=\mu^{-1}f$, so $f$
is an eigenfunction of $H$, hence lies in some $E_k\subset W$ by
\cref{isotropic:prop:HO-eigenspaces}, a contradiction.
\end{proof}

\begin{proposition}\label{isotropic:prop:hermite-schwartz}
If $u\in\cS(\R^n)$ then the coefficients in \eqref{isotropic:eq:hermite-expansion}
are rapidly decreasing, $|c_\alpha|\le C_N(1+|\alpha|)^{-N}$ for all $N$, and the
series converges absolutely in $\cS(\R^n)$. Conversely, every such series
with rapidly decreasing coefficients converges in $\cS(\R^n)$.
\end{proposition}
\begin{proof}
Since $H^Nu\in\cS$,
$(n+2|\alpha|)^N|c_\alpha|=|(u,H^Nu_\alpha)|=|(H^Nu,u_\alpha)|\le\norm{H^Nu}$. By
\eqref{isotropic:eq:ladder} and $2ix_j=C_j-A_j$, $2D_j=C_j+A_j$, each
$x^\beta D^\gamma u_\alpha$ is a combination of at most $2^{|\beta|+|\gamma|}$ functions
$u_{\alpha'}$ with coefficients bounded by $(2|\alpha|+2|\beta|+2|\gamma|)^{(|\beta|+|\gamma|)/2}$.
By the Sobolev embedding
$\sup|v|\le C\sum_{|\mu|\le n+1}\norm{D^\mu v}_{L^2}$, the Schwartz seminorms of
$u_\alpha$ are therefore bounded by $C_k(1+|\alpha|)^{M_k}$, and the series
converges absolutely in $\cS$ when the coefficients are rapidly decreasing;
its sum is $u$ since it converges to $u$ in $L^2$.
\end{proof}

For $N\in\N_0$ let $\Pi_N$ be the orthogonal projection onto
$E_0\oplus\dots\oplus E_N$, i.e.\ $K_{\Pi_N}=\sum_{|\alpha|\le N}u_\alpha(x)\overline{u_\alpha(y)}$;
it is a finite rank element of $\Psiiso{-\infty}(\R^n)$.

\begin{corollary}\label{isotropic:cor:finite-rank-dense}
For every $A\in\Psiiso{-\infty}(\R^n)$, $\Pi_NA\Pi_N\to A$ in
$\Psiiso{-\infty}(\R^n)$ as $N\to\infty$. In particular the finite rank
elements are dense in $\Psiiso{-\infty}(\R^n)$, and also in $\Lin(L^2)$ norm.
\end{corollary}
\begin{proof}
The functions $u_\alpha(x)\overline{u_\beta(y)}$, $(\alpha,\beta)\in\N_0^{2n}$, are,
up to factors $\pm1$ (note $\overline{u_\beta}=(-1)^{|\beta|}u_\beta$), the
Hermite functions of $\R^{2n}$. Hence by
\cref{isotropic:prop:hermite-schwartz} in $\R^{2n}$ the expansion
$K_A=\sum_{\alpha,\beta}c_{\alpha\beta}u_\alpha(x)\overline{u_\beta(y)}$ converges
absolutely in $\cS(\R^{2n})$; its partial sum over $|\alpha|,|\beta|\le N$ is the
kernel of $\Pi_NA\Pi_N$. The last statement follows from
\cref{isotropic:prop:compact}.
\end{proof}

\begin{remark}\label{isotropic:rem:HO-powers}
The same arguments show that for $p\in\Z$ the eigenvalues of $H^p$ are
$(n+2k)^p$, with the same eigenspaces $E_k$. Thus $H^p$ is compact if $p<0$ and
$H^p=\sum_\alpha(n+2|\alpha|)^pu_\alpha\otimes\bar u_\alpha$ is of trace class if $p<-n$.
\end{remark}

%======================================================================
\section{Isotropic Sobolev spaces}\label{isotropic:sec:sobolev}

The space $L^2(\R^n)$ is associated with the inclusion
$\cS(\R^n)\subset\cSp(\R^n)$, and operators of order $0$ are bounded on it. The
isotropic Sobolev space of order $s$ consists of the distributions mapped into
$L^2$ by all isotropic operators of order $s$. For each $s\in\R$ we fix the
elliptic operator $\Lambda_s=\Op(\ang{w}^s)\in\Psiisoph{s}(\R^n)$, and let
$P_s$ be the orthogonal projection onto its (finite-dimensional) null space
and $G_s\in\Psiisoph{-s}$ its generalized inverse
(\cref{isotropic:prop:gen-inverse}), $G_s\Lambda_s=\Id-P_s$.

\begin{definition}\label{isotropic:def:Hiso}
For $s\in\R$,
\begin{equation*}
  \Hiso{s}(\R^n)=\{u\in\cSp(\R^n);\ Au\in L^2(\R^n)\ \forall\,A\in\Psiiso{s}(\R^n)\},
\end{equation*}
with the norm $\norm{u}_{s}^2=\norm{\Lambda_su}_{L^2}^2+\norm{P_su}_{L^2}^2$.
\end{definition}

\begin{proposition}\label{isotropic:prop:Hiso}
\begin{enumerate}[label=\textup{(\roman*)}]
\item If $A\in\Psiiso{s}(\R^n)$ is elliptic then
  $\Hiso{s}(\R^n)=\{u\in\cSp;\ Au\in L^2\}$, and $\norm{Au}+\norm{\Pi_0u}$, with
  $\Pi_0$ the projection onto $\Nul(A)$, is an equivalent norm.
\item $\Hiso{s}(\R^n)$ is a Hilbert space in which $\cS(\R^n)$ is dense; $\Hiso{0}=L^2$.
\item For $k\in\N_0$,
  $\Hiso{k}(\R^n)=\{u\in L^2;\ x^\alpha D^\beta u\in L^2\ \forall\,|\alpha|+|\beta|\le k\}$
  with equivalent norm $\bigl(\sum_{|\alpha|+|\beta|\le k}\norm{x^\alpha D^\beta u}^2\bigr)^{1/2}$.
\item The $L^2$ pairing extends to a non-degenerate pairing
  $\Hiso{s}\times\Hiso{-s}\to\C$ identifying $\Hiso{-s}$ with the (anti)dual of
  $\Hiso{s}$.
\item Each $A\in\Psiiso{m}(\R^n)$ is bounded $\Hiso{s}\to\Hiso{s-m}$ for all $s$.
\item $\bigcap_s\Hiso{s}=\cS(\R^n)$, $\bigcup_s\Hiso{s}=\cSp(\R^n)$, and for $s>t$
  the inclusion $\Hiso{s}\to\Hiso{t}$ is compact.
\end{enumerate}
\end{proposition}
\begin{proof}
(i) If $A$ is elliptic of order $s$ with parametrix $B$, any $A'\in\Psiiso{s}$
can be written $A'=(A'B)A+A'(\Id-BA)$, with $A'B\in\Psiiso{0}$ bounded on $L^2$
and $A'(\Id-BA)$ residual, mapping $\cSp$ to $\cS$. So $Au\in L^2$ implies
$A'u\in L^2$. If $G$ is the generalized inverse of $A$,
$u=GAu+\Pi_0u$, so $\norm{A'u}\le\norm{A'G}\norm{Au}+\norm{A'\Pi_0}\norm{\Pi_0u}$; by
symmetry the norms are equivalent.

(ii) Let $u_k$ be Cauchy. Then $\Lambda_su_k\to f$ in $L^2$ and $P_su_k\to g$.
Since $\Pi_1\Lambda_su_k=0$ ($\Pi_1$ the projection onto $\Nul\Lambda_s^*$), also
$\Pi_1f=0$. As $G_s$ is continuous on $\cSp$, $u_k=G_s\Lambda_su_k+P_su_k\to u=G_sf+g$
in $\cSp$, and $\Lambda_su=(\Id-\Pi_1)f=f$, $P_su=g$ (note $P_sG_s=0$ by
construction). Thus $u_k\to u$ in $\Hiso{s}$. For density, given
$u\in\Hiso{s}$ choose $f_k\in\cS$ with $f_k\to\Lambda_su$ in $L^2$; then
$u_k=G_sf_k+P_su\in\cS$ and $u_k\to u$ in $\Hiso{s}$. That $\Hiso{0}=L^2$
follows from (i) with $A=\Id$.

(iii) Since $x^\alpha D^\beta\in\Psiiso{|\alpha|+|\beta|}\subset\Psiiso{k}$ for
$|\alpha|+|\beta|\le k$, elements of $\Hiso{k}$ satisfy the conditions.
Conversely the operator
$Q_k=\sum_{|\alpha|+|\beta|\le k}(x^\alpha D^\beta)^*(x^\alpha D^\beta)\in\Psiisoph{2k}$
has principal symbol $q_k=\sum_{|\alpha|+|\beta|=k}x^{2\alpha}\xi^{2\beta}$, which is
positive on $\Sphere^{2n-1}$ (if $w_i\ne0$ the term $w_i^{2k}$ occurs), so $Q_k$
is elliptic. If $B$ is a parametrix and $A\in\Psiiso{k}$, then
$Au=\sum(ABD^\beta x^\alpha)(x^\alpha D^\beta u)+A(\Id-BQ_k)u$ with
$ABD^\beta x^\alpha\in\Psiiso{0}$. The norm equivalence follows as in (i).

(iv) For $u\in\cS$, $v\in\Hiso{-s}$:
$(u,v)=(G_s\Lambda_su+P_su,v)=(\Lambda_su,G_s^*v)+(u,P_sv)$, and $G_s^*\in\Psiiso{-s}$,
so $|(u,v)|\le C\norm{u}_s\norm{v}_{-s}$; the pairing extends by density.
Conversely a continuous antilinear functional $\ell$ on $\Hiso{s}$ is, by the
Riesz theorem, $\ell(u)=\overline{\inner{w}{u}_s}$ for some $w\in\Hiso{s}$, i.e.\
$\ell(u)=\overline{(u,v)}$ with $v=\Lambda_s^*\Lambda_sw+P_sw$, and $v\in\Hiso{-s}$
since $A\Lambda_s^*\in\Psiiso{0}$ for $A\in\Psiiso{-s}$. Non-degeneracy follows
from the density of $\cS$.

(v) If $A'\in\Psiiso{s-m}$ then $A'A\in\Psiiso{s}$, so $A$ maps $\Hiso{s}$ into
$\Hiso{s-m}$, and boundedness follows from
$\norm{\Lambda_{s-m}Au}\le\norm{\Lambda_{s-m}AG_s}\norm{\Lambda_su}+\norm{\Lambda_{s-m}AP_s}\norm{P_su}$.

(vi) If $u\in\Hiso{k}$ for all $k$ then $x^\alpha D^\beta u\in L^2$ for all
$\alpha,\beta$, so $u\in\cS$ by Sobolev embedding. A tempered distribution
satisfies
$|\pair{u}{\phi}|\le C\sum_{|\alpha|+|\beta|\le k}\sup|x^\alpha D^\beta\phi|\le C'\norm{\phi}_{k+n+1}$
(Sobolev embedding again), so by (iii), (iv) it lies in $\Hiso{-k-n-1}$. For
compactness, $\Lambda_tu=\Lambda_tG_s\Lambda_su+\Lambda_tP_su$ where
$\Lambda_tG_s\in\Psiiso{t-s}$ is compact on $L^2$ (\cref{isotropic:prop:L2}) and
$\Lambda_tP_s$ has finite rank.
\end{proof}

\begin{proposition}\label{isotropic:prop:Hiso-fredholm}
An operator $A\in\Psiiso{m}(\R^n)$ is Fredholm as a map
$\Hiso{s}(\R^n)\to\Hiso{s-m}(\R^n)$ if and only if it is elliptic; then its null
space and index are those of \cref{isotropic:prop:gen-inverse}, independent of $s$.
\end{proposition}
\begin{proof}
If $A$ is elliptic, \eqref{isotropic:eq:gen-inverse} with $G$ bounded
$\Hiso{s-m}\to\Hiso{s}$ shows that $A$ is Fredholm, with null space $N_0$ and
range $\{f;\ \Pi_1f=0\}$. In particular $\Lambda_t:\Hiso{t}\to L^2$ and
$G_t:L^2\to\Hiso{t}$ are Fredholm. If $A:\Hiso{s}\to\Hiso{s-m}$ is Fredholm,
then $T=\Lambda_{s-m}AG_s\in\Psiiso{0}$ is Fredholm on $L^2$, so elliptic by
\cref{isotropic:cor:fredholm-L2}. Since $\sigma_L(T)-\ang{w}^{s-m}a\,\sigma_L(G_s)\in\Gs{-2}$
and $\ang{w}^{s-m}$ and $\sigma_L(G_s)$ are elliptic of orders $s-m$ and $-s$, $a$
is elliptic of order $m$.
\end{proof}

\begin{remark}\label{isotropic:rem:Hsc}
For the scattering calculus the natural spaces carry two orders,
$H^{m,l}_{\mathrm{sc}}(\R^n)=\ang{x}^{-l}H^m(\R^n)$, and analogous results hold
for $\Psisc{m,l}(\R^n)$ acting between them; we do not pursue this here.
\end{remark}

%======================================================================
\section{The residual group}\label{isotropic:sec:group}

\begin{definition}\label{isotropic:def:group}
The \emph{residual isotropic group} is
\begin{equation*}
  \Gres(\R^n)=\{\Id+A;\ A\in\Psiiso{-\infty}(\R^n),\ \Id+A\text{ invertible on }L^2(\R^n)\}.
\end{equation*}
\end{definition}

By \cref{isotropic:cor:invertible}, $\Gres(\R^n)$ is a group, the inverse of
$\Id+A$ being of the form $\Id+B$ with $B\in\Psiiso{-\infty}$, and $\Id+A$ is in
$\Gres$ as soon as it is injective (or surjective) on $L^2$. We identify
$\Id+\Psiiso{-\infty}(\R^n)$ with $\Psiiso{-\infty}(\R^n)\cong\cS(\R^{2n})$ and
regard $\Gres(\R^n)$ as an open subset of this Fréchet space.

\begin{proposition}\label{isotropic:prop:group}
$\Gres(\R^n)$ is open in $\Id+\Psiiso{-\infty}(\R^n)$, and multiplication and
inversion are smooth; the derivative of inversion at $g$ is
$v\mapsto-g^{-1}vg^{-1}$. The exponential map of
\cref{isotropic:prop:exp-log} takes values in $\Gres(\R^n)$ and restricts to a
homeomorphism of a neighbourhood of $0$ onto a neighbourhood of $\Id$.
\end{proposition}
\begin{proof}
Openness follows from the openness of the invertible operators in $\Lin(L^2)$
and the continuity of $\Psiiso{-\infty}\to\Lin(L^2)$
(\cref{isotropic:prop:compact}). Multiplication
$(\Id+A)(\Id+B)=\Id+A+B+AB$ is continuous and bilinear in $(A,B)$, hence
smooth. If $(\Id+A)^{-1}=\Id+B$, then as in the proof of
\cref{isotropic:prop:Id+A}, $B=-A+A^2+ABA$. Since $A\mapsto B$ is continuous
into $\Lin(L^2)$, \eqref{isotropic:eq:corner-est} shows it is continuous into
$\Psiiso{-\infty}$. From $(g+tv)^{-1}-g^{-1}=-t(g+tv)^{-1}vg^{-1}$ the
derivative of inversion is $-g^{-1}vg^{-1}$, continuous in $(g,v)$, and by
induction inversion is $\CI$. The last statement is
\cref{isotropic:prop:exp-log}.
\end{proof}

\subsection{Differential forms}

Here a map $f:U\to F$ from an open subset $U$ of a Fréchet space $E$ into a
Fréchet space $F$ is $\mathcal C^1$ if the directional derivatives
$Df(u;v)=\frac{d}{dt}f(u+tv)|_{t=0}$ exist and $(u,v)\mapsto Df(u;v)$ is
continuous on $U\times E$; $\mathcal C^k$ and $\CI$ are defined inductively.
For $E=\cS(\R^p)$, the continuous $k$-linear forms on $E$ are identified, by
the Schwartz kernel theorem, with $\cSp(\R^{kp})$, and the alternating ones with
\begin{equation*}
  \Lambda^k\cS(\R^p)=\{u\in\cSp(\R^{kp});\ u(x_{e(1)},\dots,x_{e(k)})=\sgn(e)\,u(x_1,\dots,x_k)\ \forall\,e\},
\end{equation*}
$e$ ranging over the permutations of $\{1,\dots,k\}$ and $x_i\in\R^p$. A
\emph{$k$-form} on an open set $U\subset\cS(\R^p)$ is a map
$\alpha:U\times E^k\to\C$, alternating and $k$-linear in the last $k$ arguments
(so $\alpha(u)\in\Lambda^k\cS(\R^p)$ for each $u$), which is smooth on
$U\times E^k$. The exterior derivative is defined by the usual formula for
constant vector fields,
\begin{equation*}
  d\alpha(u;v_0,\dots,v_k)=\sum_{i=0}^k(-1)^iD_u\alpha(u;v_0,\dots,\widehat{v_i},\dots,v_k)[v_i],
\end{equation*}
so that $df(u;v)=Df(u;v)$ for functions; the symmetry of second derivatives
gives $d^2=0$, and $d(\alpha\wedge\beta)=d\alpha\wedge\beta+(-1)^{\deg\alpha}\alpha\wedge d\beta$.
Applying this with $U=\Gres(\R^n)$, $E=\Psiiso{-\infty}(\R^n)\cong\cS(\R^{2n})$,
we obtain the de~Rham complex of $\Gres(\R^n)$; for example the
Maurer--Cartan form $g^{-1}dg$, $(g;v)\mapsto g^{-1}v$, is a
$\Psiiso{-\infty}(\R^n)$-valued $1$-form. These forms are used in
\cref{ch:determinants}.

%======================================================================
\section{Symplectic invariance and the Weyl quantization}\label{isotropic:sec:symplectic}

The composition law of the isotropic calculus has a more invariant
formulation, which exhibits its invariance under linear symplectic maps.

\subsection{Symplectic vector spaces}

A \emph{symplectic vector space} is a real vector space $W$ with an
antisymmetric non-degenerate bilinear form $\omega$:
\begin{equation*}
  \omega(w_1,w_2)+\omega(w_2,w_1)=0,\qquad\omega(v,w)=0\ \forall\,w\in W\implies v=0 .
\end{equation*}
A subspace $V$ is \emph{isotropic} if $\omega$ vanishes on $V\times V$, and
\emph{Lagrangian} if moreover $2\dim V=\dim W$.

\begin{lemma}\label{isotropic:lem:lagrangian}
Every symplectic vector space $W$ has a Lagrangian subspace, and for every
Lagrangian subspace $U_1$ there is a Lagrangian subspace $U_2$ with
$W=U_1\oplus U_2$.
\end{lemma}
\begin{proof}
Every one-dimensional subspace is isotropic, by antisymmetry. Let $V$ be an
isotropic subspace of maximal dimension and $U$ a complement of $V$. The
pairing $\omega:V\times U\to\R$ is non-degenerate: if $u\in U$ and
$\omega(v,u)=0$ for all $v\in V$ then $V+\R u$ is isotropic, so $u=0$ by
maximality; if $v\in V$ and $\omega(v,u)=0$ for all $u\in U$ then, since $\omega$
vanishes on $V$, $\omega(v,\cdot)=0$ and $v=0$. Hence $\dim U=\dim V$ and $V$ is
Lagrangian.

Now let $U_1$ be Lagrangian, $\dim W=2n$, and let $V$ be an isotropic subspace
with $V\cap U_1=\{0\}$, of maximal dimension among such subspaces. Suppose
$\dim V<n$. The map $V\to U_1'$, $v\mapsto\omega(v,\cdot)|_{U_1}$, is injective,
for if $\omega(v,U_1)=0$ then $U_1+\R v$ is isotropic, so $v\in U_1$ and $v=0$.
Hence $U_1\to V'$, $u\mapsto\omega(u,\cdot)|_V$, is surjective. Choose
$w\notin V\oplus U_1$, possible since $\dim V+\dim U_1<2n$, and $u_1\in U_1$
with $\omega(u_1,v)=\omega(w,v)$ for all $v\in V$. Then $V\oplus\R(w-u_1)$ is
isotropic and meets $U_1$ only in $0$ (as $w\notin V+U_1$), contradicting
maximality. Thus $\dim V=n$ and $U_2=V$ is a Lagrangian complement.
\end{proof}

Given a Lagrangian decomposition $W=U_1\oplus U_2$, the pairing
$\omega:U_2\times U_1\to\R$ is non-degenerate. Choosing a basis $e_1,\dots,e_n$ of
$U_1$ and the basis $f_1,\dots,f_n$ of $U_2$ with $\omega(f_i,e_j)=\delta_{ij}$,
and writing $w=\sum x_ie_i+\sum\xi_if_i$,
\begin{equation}\label{isotropic:eq:darboux}
  \omega(w,w')=\sum_{i=1}^n(\xi_ix_i'-x_i\xi_i'),\qquad\text{i.e.}\quad\omega=\sum_{i=1}^nd\xi_i\wedge dx_i .
\end{equation}
These are \emph{Darboux coordinates}. The \emph{Liouville measure}
$\nu=|\omega^n/n!|$ is, in Darboux coordinates, Lebesgue measure $dx\,d\xi$,
and is invariant under the symplectic group $\mathrm{Sp}(W)$ of linear maps
preserving $\omega$.

\subsection{Weyl quantization}

For $a\in\Gs{m}(\R^{2n})$ the \emph{Weyl quantization} of $a$ is
\begin{equation*}
  \OpW(a)u(x)=\int\!\!\int e^{i(x-y)\cdot\xi}a\bigl(\tfrac{x+y}2,\xi\bigr)u(y)\,dy\,\dbar\xi ,
\end{equation*}
i.e.\ $\OpW(a)=\Op(a(\frac{x+y}2,\xi))$. Its kernel $K$ and $a$ are related by
\begin{equation}\label{isotropic:eq:weyl-kernel}
  K\bigl(s+\tfrac t2,s-\tfrac t2\bigr)=\int e^{it\cdot\xi}a(s,\xi)\,\dbar\xi,\qquad
  a(s,\xi)=\int e^{-it\cdot\xi}K\bigl(s+\tfrac t2,s-\tfrac t2\bigr)\,dt .
\end{equation}

\begin{proposition}\label{isotropic:prop:weyl}
For $a\in\Gs{m}(\R^{2n})$, $\OpW(a)\in\Psiiso{m}(\R^n)$ with
\begin{equation*}
  \sigma_L(\OpW(a))\sim\sum_\alpha\frac{(i/2)^{|\alpha|}}{\alpha!}D_\xi^\alpha D_x^\alpha a ,
\end{equation*}
and $\OpW(a)^*=\OpW(\bar a)$. Conversely every $A\in\Psiiso{m}(\R^n)$ is
$\OpW(a)$ for a unique $a\in\Gs{m}(\R^{2n})$, the \emph{Weyl symbol} of $A$;
$A\in\Psiisoph{m}$ if and only if $a\in\Gph{m}$, and $A\in\Psiiso{-\infty}$
if and only if $a\in\cS(\R^{2n})$.
\end{proposition}
\begin{proof}
The first statement is \cref{isotropic:prop:reduction} applied to the
amplitude $a(\frac{x+y}2,\xi)\in\Ampl{m}$
(\cref{isotropic:lem:ampl-examples}(ii) with $\lambda=\frac12$), since
$D_y^\alpha[a(\frac{x+y}2,\xi)]|_{y=x}=2^{-|\alpha|}D_x^\alpha a(x,\xi)$. The
adjoint formula is immediate from the kernel. By
\eqref{isotropic:eq:weyl-kernel}, $a$ is determined by $K$, and
$a\in\cS(\R^{2n})$ if and only if $K\in\cS(\R^{2n})$. Given $A$ with left symbol
$b$, choose $a'\in\Gs{m}$ with $a'\sim\sum_\alpha\frac{(-i/2)^{|\alpha|}}{\alpha!}D_\xi^\alpha D_x^\alpha b$;
by the formal identity $e^{\frac i2D_x\cdot D_\xi}e^{-\frac i2D_x\cdot D_\xi}=1$,
$A-\OpW(a')\in\Psiiso{-\infty}=\OpW(\cS)$, so $A=\OpW(a)$ with $a-a'\in\cS$. The
expansions preserve polyhomogeneity.
\end{proof}

\begin{proposition}\label{isotropic:prop:weyl-product}
If $a,b\in\cS(\R^{2n})$ then $\OpW(a)\OpW(b)=\OpW(a\#_{\mathrm W}b)$ where
\begin{equation}\label{isotropic:eq:weyl-product}
  a\#_{\mathrm W}b(w)=\pi^{-2n}\int\!\!\int_{\R^{2n}\times\R^{2n}}e^{-2i\omega(w_1,w_2)}a(w+w_1)b(w+w_2)\,dw_1\,dw_2,
\end{equation}
$\omega$ being the standard symplectic form \eqref{isotropic:eq:darboux} on $\R^{2n}$.
\end{proposition}
\begin{proof}
Let $A=\OpW(a)$, $B=\OpW(b)$ with kernels $K_A,K_B$, and let $c$ be the Weyl
symbol of $AB$. By \eqref{isotropic:eq:weyl-kernel},
\begin{align*}
  c(s,\zeta)&=\int e^{-it\cdot\zeta}\int K_A\bigl(s+\tfrac t2,z\bigr)K_B\bigl(z,s-\tfrac t2\bigr)\,dz\,dt\\
  &=(2\pi)^{-2n}\int e^{i\Phi}a\bigl(\tfrac s2+\tfrac t4+\tfrac z2,\xi\bigr)b\bigl(\tfrac z2+\tfrac s2-\tfrac t4,\eta\bigr)\,dt\,dz\,d\xi\,d\eta,\\
  \Phi&=-t\cdot\zeta+\bigl(s+\tfrac t2-z\bigr)\cdot\xi+\bigl(z-s+\tfrac t2\bigr)\cdot\eta,
\end{align*}
all integrals converging absolutely. Change variables to
$X=\frac z2+\frac t4-\frac s2$, $Y=\frac z2-\frac t4-\frac s2$, $\Xi=\xi-\zeta$,
$\mathrm H=\eta-\zeta$. Then $z-s=X+Y$, $t=2(X-Y)$, $dt\,dz=4^n\,dX\,dY$ and
$\Phi=2(X\cdot\mathrm H-Y\cdot\Xi)$, so
\begin{equation*}
  c(s,\zeta)=\pi^{-2n}\int e^{2i(X\cdot\mathrm H-Y\cdot\Xi)}a(s+X,\zeta+\Xi)b(s+Y,\zeta+\mathrm H)\,dX\,d\Xi\,dY\,d\mathrm H .
\end{equation*}
With $w=(s,\zeta)$, $w_1=(X,\Xi)$, $w_2=(Y,\mathrm H)$ we have
$\omega(w_1,w_2)=\Xi\cdot Y-X\cdot\mathrm H$, which gives \eqref{isotropic:eq:weyl-product}.
\end{proof}

\begin{theorem}\label{isotropic:thm:symplectic}
Let $(W,\omega)$ be a symplectic vector space of dimension $2n$ with Liouville
measure $\nu$. Then
\begin{equation}\label{isotropic:eq:symplectic-product}
  a\#b(w)=\pi^{-2n}\int_{W\times W}e^{-2i\omega(w_1,w_2)}a(w+w_1)b(w+w_2)\,d\nu(w_1)\,d\nu(w_2)
\end{equation}
defines an associative product on $\cS(W)$. For any Lagrangian decomposition
$W=U_1\oplus U_2$ and Darboux coordinates as above, $(\cS(W),\#)$ is isomorphic,
via the Weyl quantization, to $(\Psiiso{-\infty}(\R^n),\circ)$. The product
extends uniquely to continuous bilinear maps $\Gs{m}(W)\times\Gs{m'}(W)\to\Gs{m+m'}(W)$,
which in any Darboux coordinates correspond to composition in $\Psiiso{\infty}(\R^n)$,
and for every $F\in\mathrm{Sp}(W)$
\begin{equation}\label{isotropic:eq:symplectic-inv}
  F^*(a\#b)=(F^*a)\#(F^*b).
\end{equation}
\end{theorem}
\begin{proof}
In Darboux coordinates \eqref{isotropic:eq:symplectic-product} becomes
\eqref{isotropic:eq:weyl-product}, so the isomorphism (and hence associativity)
follows from \cref{isotropic:prop:weyl-product,isotropic:prop:weyl}. For
$F\in\mathrm{Sp}(W)$, substituting $v_i=Fw_i$ in
\eqref{isotropic:eq:symplectic-product}, which preserves $\omega$ and $\nu$, gives
\eqref{isotropic:eq:symplectic-inv} for $a,b\in\cS(W)$. In Darboux coordinates
the Weyl product of symbols in $\Gs{m}$, $\Gs{m'}$ is continuous into
$\Gs{m+m'}$: it is the composite of the continuous maps $a\mapsto\sigma_L(\OpW(a))$,
its inverse (both continuous by the closed graph theorem, as in
\cref{isotropic:thm:algebra}(iv)) and the product of
\cref{isotropic:thm:algebra}(iv). Since $\cS$ is dense in $\Gs{m}$ for the
topology of $\Gs{m+\eps}$ (\cref{isotropic:lem:Gamma-basic}(iv)) and $\eps>0$ is
arbitrary, the extension is unique, and \eqref{isotropic:eq:symplectic-inv}
extends by continuity because $F^*$ is continuous on each $\Gs{m}(W)$. In
particular, since two Darboux coordinate systems differ by an element of
$\mathrm{Sp}(W)$, the extended product does not depend on the choice.
\end{proof}

\begin{remark}\label{isotropic:rem:metaplectic}
The Fourier transform is an instance of this invariance. A computation with
the operators $\OpW(e^{i(p\cdot x+q\cdot\xi)})=e^{i(p\cdot x+q\cdot D)}$ (or a
comparison with \cref{isotropic:prop:fourier}) shows that the Weyl symbol of
$\cF A\cF^{-1}$ is $a\circ J$, where $J(x,\xi)=(-\xi,x)$ is symplectic. More
generally each $F\in\mathrm{Sp}(\R^{2n})$ is implemented, in the sense that
$\OpW(F^*a)=U_F^{-1}\OpW(a)U_F$, by a unitary operator $U_F$ on $L^2(\R^n)$,
determined up to a constant factor of modulus one; these operators generate the
metaplectic group, see \cite[Section~18.5]{Hormander-III}.
\end{remark}

%======================================================================
\section{Complex order}\label{isotropic:sec:complex}

The identification $\Gph{0}(\R^{2n})=\rc^*\CI(\Hsph{2n})$ of
\cref{isotropic:prop:rc} allows us to define isotropic operators of complex
order. For $z\in\C$ set
\begin{equation*}
  \Gph{z}(\R^N)=\ang{w}^z\,\Gph{0}(\R^N)=\rc^*\bigl(\rho^{-z}\CI(\Hsph{N})\bigr),\qquad
  \Psiiso{z}(\R^n)=\Op\bigl(\Gph{z}(\R^{2n})\bigr).
\end{equation*}
For real $z$ this agrees with $\Gph{z}$ as defined before, since $\ang{w}^z\in\Gph{z}$.
Since $|\pa^\alpha\ang{w}^z|\le C_\alpha\ang{w}^{\operatorname{Re}z-|\alpha|}$,
$\Gph{z}\subset\Gs{\operatorname{Re}z}$ and $\Psiiso{z}\subset\Psiiso{\operatorname{Re}z}$. The
definition does not depend on the choice of boundary defining function: if
$\rho'=f\rho$ with $0<f\in\CI(\Hsph{N})$ then $f^{z}=e^{z\log f}$ is smooth, so
$\rho'^{-z}\CI=\rho^{-z}\CI$.

\begin{proposition}\label{isotropic:prop:complex-order}
For $z,z'\in\C$, $\Psiiso{z}(\R^n)\circ\Psiiso{z'}(\R^n)\subset\Psiiso{z+z'}(\R^n)$
and $\Psiiso{z}(\R^n)^*=\Psiiso{\bar z}(\R^n)$.
\end{proposition}
\begin{proof}
Differentiation maps $\ang{w}^z\Gph{0}$ into $\ang{w}^z\Gph{-1}$, since
$\pa_{w_j}(\ang{w}^zf)=\ang{w}^z(zw_j\ang{w}^{-2}f+\pa_{w_j}f)$. Hence the terms
of \eqref{isotropic:eq:composition} satisfy
$\frac1{\alpha!}\pa_\xi^\alpha a\,D_x^\alpha b\in\ang{w}^{z+z'}\Gph{-2|\alpha|}$. If
$c=a\#b$ then $c\in\Gs{\mu}$, $\mu=\operatorname{Re}(z+z')$, and the remainders of the
expansion lie in $\Gs{\mu-2N}$; multiplying by $\ang{w}^{-(z+z')}\in\Gs{-\mu}$,
$\ang{w}^{-(z+z')}c$ has an asymptotic expansion with terms in $\Gph{-2j}$ and
remainders in $\Gs{-2N}$, so it lies in $\Gph{0}$. Adjoints are treated in the
same way.
\end{proof}

Even more useful are holomorphic families. A map $\Omega\to\CI(\Hsph{N})$, with
$\Omega\subset\C$ open, is \emph{holomorphic} if it is given by an element of
$\CI(\Omega\times\Hsph{N})$ satisfying the Cauchy--Riemann equation in the first
variable. If $h:\Omega\to\C$ is holomorphic, a family $A(z)\in\Psiiso{h(z)}(\R^n)$,
$z\in\Omega$, is a \emph{holomorphic family of order $h$} if
$\sigma_L(A(z))=\ang{w}^{h(z)}\rc^*F(z)$ with $F:\Omega\to\CI(\Hsph{2n})$
holomorphic.

\begin{proposition}\label{isotropic:prop:holomorphic}
If $h,g$ are holomorphic on an open set $\Omega\subset\C$ and $A(z)$, $B(z)$ are
holomorphic families of isotropic operators of orders $h(z)$ and $g(z)$, then
$A(z)\circ B(z)$ is a holomorphic family of order $h(z)+g(z)$.
\end{proposition}
\begin{proof}[Sketch of proof]
It suffices to work on an open $\Omega'$ with compact closure in $\Omega$, on
which $\operatorname{Re}h,\operatorname{Re}g<M$ and $\operatorname{Re}(h+g)>-M$. The map
$z\mapsto\ang{w}^{h(z)}$ is holomorphic into $\Gs{M}$ (its derivative
$h'(z)\log\ang{w}\,\ang{w}^{h(z)}$ lies in $\Gs{M}$, as $\operatorname{Re}h<M$
strictly), and $\rc^*:\CI(\Hsph{2n})\to\Gs{0}$ is continuous linear, so
$a(z)=\sigma_L(A(z))$ and $b(z)$ are holomorphic $\Gs{M}$-valued. By
\cref{isotropic:thm:algebra}(iv), $c(z)=a(z)\#b(z)$ is holomorphic into
$\Gs{2M}$ and
\begin{equation*}
  c(z)=\sum_{|\alpha|<N}\frac1{\alpha!}\pa_\xi^\alpha a(z)\,D_x^\alpha b(z)+R_N(z),
\end{equation*}
with $R_N$ holomorphic into $\Gs{2M-2N}$. Using \eqref{isotropic:eq:b-fields},
$\pa_{w_j}(\ang{w}^h\rc^*F)=\ang{w}^{h-1}\rc^*\bigl(hX_jF+V_jF\bigr)$ with
$X_j=w_j\ang{w}^{-1}$ and $V_j$ smooth on $\Hsph{2n}$, so
$\ang{w}^{-h-g}$ times each term of the sum is $\rc^*$ of a holomorphic family
in $\CI(\Hsph{2n})$, while $\ang{w}^{-h-g}R_N(z)$ is holomorphic into
$\Gs{3M-2N}$, hence, by \cref{isotropic:prop:rc}(iii), into
$\mathcal C^{k}(\Hsph{2n})$ for $2N>3M+k$. Thus $\rho^{h+g}c(z)$ (pushed forward
by $\rc$) is holomorphic into $\mathcal C^k(\Hsph{2n})$ for every $k$, which
means that it is given by an element of $\CI(\Omega'\times\Hsph{2n})$
holomorphic in $z$.
\end{proof}

%======================================================================
\section{Exercises}\label{isotropic:sec:exercises}

\begin{exercise}\label{isotropic:ex:sobolev-integer}
For $k\in\N$ define $\Hiso{-k}(\R^n)$ as the space of $u\in\cSp(\R^n)$ which can
be written $u=\sum_{|\alpha|+|\beta|\le k}x^\alpha D^\beta u_{\alpha\beta}$ with
$u_{\alpha\beta}\in L^2(\R^n)$. Show directly that each $A\in\Psiiso{p}(\R^n)$,
$p\in\Z$, maps $\Hiso{k}$ to $\Hiso{k-p}$ for all $k\in\Z$ (with $\Hiso{k}$,
$k\ge0$, as in \cref{isotropic:prop:Hiso}(iii)), and that these definitions
agree with \cref{isotropic:def:Hiso}. Deduce that if $A\in\Psiiso{m}(\R^n)$ is
elliptic, $u\in\cSp(\R^n)$ and $Au\in\Hiso{m'}(\R^n)$, then
$u\in\Hiso{m'+m}(\R^n)$.
\end{exercise}

\begin{exercise}\label{isotropic:ex:sobolev-compact}
Show that for $\eps>0$,
$\Hiso{\eps}(\R^n)\subset\ang{x}^{-\eps}L^2(\R^n)\cap H^\eps(\R^n)$, and that the
inclusion is strict. [Hint: compose with $\ang{D}^\eps$ and $\ang{x}^\eps$,
using the scattering calculus.] Deduce, without using
\cref{isotropic:prop:L2}, that $\Hiso{\eps}(\R^n)\hookrightarrow L^2(\R^n)$ is
compact, and hence that each element of $\Psiiso{-\eps}(\R^n)$ is compact on
$L^2(\R^n)$.
\end{exercise}

\begin{exercise}\label{isotropic:ex:finite-rank-parametrix}
Show that if $E\in\Psiiso{-\infty}(\R^n)$ there exists
$F\in\Psiiso{-\infty}(\R^n)$ such that $(\Id+E)(\Id+F)=\Id+G$ with
$G\in\Psiiso{-\infty}(\R^n)$ of finite rank. Deduce that an elliptic
$A\in\Psiiso{m}(\R^n)$ has a parametrix $B\in\Psiiso{-m}(\R^n)$ up to finite
rank errors, i.e.\ $AB-\Id$ and $BA-\Id$ are finite rank elements of
$\Psiiso{-\infty}(\R^n)$, and hence that $A:\Hiso{M}(\R^n)\to\Hiso{M-m}(\R^n)$ is
Fredholm for every $M$. (Recall that a bounded operator between Hilbert spaces
is Fredholm if it has finite-dimensional null space and closed range of
finite codimension.)
\end{exercise}

\begin{exercise}\label{isotropic:ex:index-zero}
Suppose $A\in\Psiiso{m}(\R^n)$ is elliptic and has index zero as an operator on
$\cSp(\R^n)$. Show that there is a finite rank $E\in\Psiiso{-\infty}(\R^n)$
such that $A+E$ is an isomorphism of $\cSp(\R^n)$, of $\cS(\R^n)$ and of
$\Hiso{s}(\R^n)\to\Hiso{s-m}(\R^n)$ for every $s$, with inverse in $\Psiiso{-m}(\R^n)$.
\end{exercise}

\begin{exercise}\label{isotropic:ex:koszul}
[Koszul complex] Let $\Lambda^k=\Lambda^k\R^n$ with basis $dx^I$, $I$ strictly
increasing $k$-tuples, declared orthonormal, and let $\cSp(\R^n;\Lambda^k)$ be
the space of $u=\sum_Iu_Idx^I$, $u_I\in\cSp(\R^n)$. Using the annihilation
operators define
\begin{equation*}
  D:\cSp(\R^n;\Lambda^k)\to\cSp(\R^n;\Lambda^{k+1}),\qquad Du=\sum_{j=1}^n\sum_IA_ju_I\,dx_j\wedge dx^I .
\end{equation*}
\begin{enumerate}[label=(\alph*)]
\item Show that $D^2=0$ and that the formal adjoint is
  $D^*u=\sum_j\sum_IC_ju_I\,\iota_jdx^I$, where $\iota_j$, contraction with $e_j$,
  is the adjoint of $dx_j\wedge$.
\item Show that $D+D^*\in\Psiiso{1}(\R^n;\Lambda^*)$ is elliptic, and that
  $(D+D^*)^2=DD^*+D^*D$ acts on $\cSp(\R^n;\Lambda^k)$ as $H-n+2k$.
\item Deduce that the null space of $D+D^*$ on $\cSp(\R^n;\Lambda^*)$ is
  $\C\exp(-|x|^2/2)\subset\cSp(\R^n;\Lambda^0)$, that
  $\{u\in\cSp(\R^n;\Lambda^0);\ Du=0\}=\C\exp(-|x|^2/2)$, and that for $k\ge1$
  $\{u\in\cSp(\R^n;\Lambda^k);\ Du=0\}=D\,\cSp(\R^n;\Lambda^{k-1})$.
\item Conclude that $D+D^*$, as an element of
  $\Psiiso{1}(\R^n;\Lambda^{\mathrm{even}},\Lambda^{\mathrm{odd}})$, is elliptic of index $1$.
\end{enumerate}
\end{exercise}

\begin{exercise}\label{isotropic:ex:WF-operator}
[Isotropic operator wavefront set] For $a\in\Gs{m}(\R^{2n})$ define the
\emph{cone support} $\operatorname{cone\,supp}(a)\subset\R^{2n}\setminus\{0\}$ as
the complement of the set of $\bar w\ne0$ having a conic neighbourhood
$\Gamma$ such that $\sup_{w\in\Gamma,|w|\ge1}\ang{w}^N|\pa^\alpha a(w)|<\infty$ for
all $N$ and $\alpha$. For $A\in\Psiiso{m}(\R^n)$ set
$\WFp_{\mathrm{iso}}(A)=\operatorname{cone\,supp}(\sigma_L(A))$. Show that
\begin{equation*}
  \WFp_{\mathrm{iso}}(A+B)\subset\WFp_{\mathrm{iso}}(A)\cup\WFp_{\mathrm{iso}}(B),\quad
  \WFp_{\mathrm{iso}}(AB)\subset\WFp_{\mathrm{iso}}(A)\cap\WFp_{\mathrm{iso}}(B),
\end{equation*}
$\WFp_{\mathrm{iso}}(A^*)=\WFp_{\mathrm{iso}}(A)$, that $\WFp_{\mathrm{iso}}(A)$ is the
same when defined using the Weyl symbol, and that
$\WFp_{\mathrm{iso}}(A)=\emptyset$ if and only if $A\in\Psiiso{-\infty}(\R^n)$.
\end{exercise}

\begin{exercise}\label{isotropic:ex:partition}
[Isotropic partition of unity] Let $U_i$, $i=1,\dots,N$, be an open cover of
$\Sphere^{2n-1}$ and $\tilde U_i=\{w\ne0;\ w/|w|\in U_i\}$. Show that there are
$A_i\in\Psiisoph{0}(\R^n)$ with $\WFp_{\mathrm{iso}}(A_i)\subset\tilde U_i$ and
$\Id-\sum_iA_i\in\Psiiso{-\infty}(\R^n)$.
\end{exercise}

\begin{exercise}\label{isotropic:ex:WF-distribution}
[Isotropic wavefront set] For $u\in\cSp(\R^n)$ let
\begin{equation*}
  \WFiso(u)=\bigcap\bigl\{\Char_{\mathrm{iso}}(A);\ A\in\Psiisoph{0}(\R^n),\ Au\in\cS(\R^n)\bigr\},
\end{equation*}
where $\Char_{\mathrm{iso}}(A)$ is the cone over the zero set of $\sigma_0(A)$ in
$\Sphere^{2n-1}$. Show that $\WFiso(u)$ is a closed cone, that
$\WFiso(u)=\emptyset$ if and only if $u\in\cS(\R^n)$, that
$\WFiso(Au)\subset\WFp_{\mathrm{iso}}(A)\cap\WFiso(u)$ for $A\in\Psiiso{m}(\R^n)$,
and that $\WFiso(\hat u)=\{(\xi,-x);\ (x,\xi)\in\WFiso(u)\}$.
\end{exercise}

\begin{exercise}\label{isotropic:ex:sc-fourier}
Verify in detail that $\cF\,\Psisc{m,l}(\R^n)\,\cF^{-1}=\Psisc{l,m}(\R^n)$, and that
an elliptic element of $\Psisc{m,l}(\R^n)$ with $l\ne0$ cannot be an element of
$\Psiiso{\infty}(\R^n)$ unless it is elliptic there of order $m=l$.
\end{exercise}

\section*{Notes}

This chapter follows Sections 3.1--3.13 of the original notes, reorganised.
The general (non-classical) isotropic class $\Gs{m}$ is the symbol class
introduced by Shubin \cite{Shubin}, where much of the material here (in
particular the harmonic oscillator and the isotropic Sobolev spaces) can be
found; the Weyl calculus and its symplectic invariance are treated in
\cite[Chapter~18.5]{Hormander-III}. The scattering calculus on $\R^n$ is the
model for Melrose's scattering calculus on asymptotically Euclidean manifolds
(R.~B.~Melrose, \emph{Spectral and scattering theory for the Laplacian on
asymptotically Euclidean spaces}, 1994); see also the work of H.~O.~Cordes and
C.~Parenti on SG-calculi.

The main departures from the original are as follows. The reduction theorem is
proved with a new class of amplitudes (\cref{isotropic:def:ampl}) which
contains the products $a(x,\xi)b(y,\xi)$ and the Weyl amplitudes, and the
isotropic decay of the remainder is obtained by intersecting weighted
traditional symbol spaces; this replaces the original argument with a
two-region decomposition of $\R^{3n}$. The generalized inverse of an elliptic
operator is constructed on $\cSp$ by a finite-dimensional argument, the
essential norm estimate is proved with coherent states (which also give the
necessity of ellipticity for the Fredholm property), the harmonic oscillator
section uses the generalized inverse $H^{-1}$ to prove completeness of the
Hermite functions, and the Sobolev spaces of all real orders are treated
uniformly. Problems on traces and determinants from the original list are in
\cref{ch:determinants}.
